% =====================================================================
%  ch10_disall.tex —— DIS 的全阶动量空间证明、重夸克与核靶（作者 W12）
%  所有不在 VERIFIED_FORMULAS_v2 中的公式与非平凡步骤均由
%  writing/work/W12/ 下的脚本重新推导（日志见 writing/work/W12/logs/）。
% =====================================================================
\chapter{DIS 的全阶动量空间证明、重夸克与核靶}\label{ch:disall}

深度非弹散射（deep inelastic scattering，DIS）是最简单的需要部分子分布的硬过程：初态只有一个强子，末态被完全求和。
\cref{ch:dis}用算符乘积展开证明了结构函数\emph{矩}的因子化（\cref{thm:dis-moments}，以 OPE 为前提而严格），\cref{ch:pdf}把它翻译为 $x$ 空间中扭度 2 部分的因子化（\cref{thm:pdf-dis-xspace}，同一等级）；
两者都只在分布意义下成立，对固定 $x$ 的逐点余项不提供控制。本章用\cref{thm:method-template}（全阶论证，[A]）给出 DIS 共线因子化的动量空间证明，
结论是固定 $0<x<1$ 处的逐点陈述（\cref{thm:disall-factorization}，\tagA）。此外推导 $\order{\as}$ 系数函数与 DIS 方案，并讨论重夸克质量与核靶；各陈述的等级与前提见下框，全阶论证的瓶颈逐一以假设的标签列出。

\begin{tcolorbox}[enhanced,breakable,title=本章证明了什么／没有证明什么,fonttitle=\bfseries,
  colback=gray!3!white,colframe=gray!55!black,boxrule=0.6pt,arc=2pt,
  left=5pt,right=5pt,top=3pt,bottom=3pt,before skip=6pt,after skip=6pt]
\textbf{证明了：}\cref{prop:disall-nlo,prop:disall-scale,prop:disall-dis-scheme}（\tagR，$R_{\mathrm{PT}}$）；\cref{cor:disall-consistency}（\tagR，$R_{\mathrm{OPE}}$）；
\cref{prop:disall-two-routes}（\tagR，$R_{\mathrm{cond}}$，前提\cref{asm:appN-collins,asm:appN-contour,asm:appN-envelope,asm:appN-separation,asm:appN-cutsum}）；
\cref{prop:disall-factorized-subtraction}（\tagR，$R_{\mathrm{cond}}$，前提\cref{asm:appN-ward-product,asm:appN-uv}）；
\cref{thm:disall-renorm-compat}（\tagR，$R_{\mathrm{cond}}$，前提\cref{asm:appN-multren,asm:appN-ir-regulator,asm:appN-contour}，内容等级 [C]）。\par
\textbf{给出全阶论证：}\cref{thm:disall-factorization}（\tagA；瓶颈\cref{asm:appN-collins,asm:appN-ward-product,asm:appN-cutsum}）；
\cref{lem:disall-fsi,prop:disall-regions,prop:disall-h2,prop:disall-h3h4,prop:disall-h5,thm:disall-heavy,prop:disall-nuclear}（\tagA）。\par
\textbf{没有证明：}$x\to1$ 处余项的一致性（阈值区）；非定域算符的 BRST 分类（具名缺口\cref{asm:appN-brst-nonlocal}，经由\cref{lem:method-ward-multi}进入）；
积分范数形式的减除估计（\cref{asm:appN-collins}）；与重整化相容的部分子靶红外正规化（\cref{asm:appN-ir-regulator}）。
\end{tcolorbox}

\paragraph{本章的思路。}
被证明因子化的首先是非物理区 $|\omega|<1$ 中的前向 Compton 振幅 $T^{\mu\nu}$（\cref{def:dis-compton}），结构函数则是它沿 $\omega\ge1$ 的不连续部分；
余项在实线段上的估计不能延拓为不连续部分的估计，所以要区分“矩”与“$x$ 空间”两条途径（\cref{sec:disall-regions}），这是 DIS 证明的特殊之处。
记号沿用\cref{ch:method}：区域 $R$、近似子 $T_R$、减除项 $C_R\Gamma$（\cref{def:method-subtraction}）、条件 (CS1)–(CS3)（\cref{def:method-collins-conditions}）、积分范数意义下的幂次压低 $\mathfrak S(U)$（\cref{def:method-power-suppressed}）；
对部分子靶，$\lambda$ 是红外正规化参数与 $Q$ 之比（例如靶的离壳度 $m/Q$），$L\equiv\ln(1/\lambda)$。

% ---------------------------------------------------------------------
\section{Breit 系与领头区域；两条途径}\label{sec:disall-regions}

先固定参考系与 $D$ 维投影算符（$W^{\mu\nu}$、$T^{\mu\nu}$ 见\cref{def:dis-kinematics,def:dis-compton}）。

\begin{definition}[Breit 系与投影算符]\label{def:disall-breit}
设靶强子动量 $P$（质量 $M$）沿 $+z$ 方向，光锥分量记作 $(v^+,v^-,\mathbf v_T)$。取
\begin{equation}\label{eq:disall-breit}
P=\Big(P^+,\frac{M^2}{2P^+},\mathbf 0_T\Big),\qquad q=\Big(-x_NP^+,\frac{Q^2}{2x_NP^+},\mathbf 0_T\Big),\qquad x_N\equiv\frac{2x}{1+\rho},\quad \rho\equiv\sqrt{1+\frac{4x^2M^2}{Q^2}},
\end{equation}
则 $q^2=-Q^2$ 与 $2P\cdot q=Q^2/x$ 精确成立，$x_N=x-x^3M^2/Q^2+\order{M^4/Q^4}$（Nachtmann 变量）。$P^+$ 可任意选取（沿 $z$ 轴的推动）；$P^+=Q/(\sqrt2\,x_N)$ 时 $q=(0,\mathbf 0_T,-Q)$（笛卡儿分量），即 Breit（“砖墙”）系本义。
在 $D=4-2\varepsilon$ 维中以同样的张量分解定义 $W^{\mu\nu}$，则对无质量靶
\begin{equation}\label{eq:disall-proj}
\frac{F_L}{x}=\frac{8x^2}{Q^2}\,P_\mu P_\nu W^{\mu\nu},\qquad
\frac{F_2}{x}=\frac1{1-\varepsilon}\Big[-g_{\mu\nu}+(3-2\varepsilon)\frac{4x^2}{Q^2}P_\mu P_\nu\Big]W^{\mu\nu};   % VF-DIS-2
\end{equation}
对有质量靶，以 $F_L\equiv\rho^2F_2-2xF_1$ 定义纵向结构函数，
\begin{equation}\label{eq:disall-proj-massive}
\frac{F_L}{x}=\frac{8x^2}{Q^2\rho^2}\,P_\mu P_\nu W^{\mu\nu},\qquad
\frac{\rho^2F_2}{x}=\frac{2(-g_{\mu\nu}W^{\mu\nu})+(D-1)F_L/x}{D-2}.   % VF-DIS-2
\end{equation}
对动量为 $p$ 的部分子靶，以 $p$ 代替 $P$、以 $z\equiv Q^2/(2p\cdot q)$ 代替 $x$，得到部分子结构函数 $\hat F_{2,L}$。
\uses{def:dis-kinematics}
\end{definition}

\noindent
投影算符来自直接缩并：记 $\tilde P=P-(P\cdot q/q^2)q$，$\tilde P^2=Q^2\rho^2/(4x^2)$，则 $-g_{\mu\nu}W^{\mu\nu}=(D-1)F_1-\tilde P^2F_2/(P\cdot q)$，
$P_\mu P_\nu W^{\mu\nu}=\frac{Q^2\rho^2}{8x^3}(\rho^2F_2-2xF_1)$。常用的 $q^+=-xP^+$ 给出 $2P\cdot q=Q^2/x-xM^2$，只在领头幂次正确。对无质量靶，被击夸克的虚度为
\begin{equation}\label{eq:disall-current-virtuality}
(\xi P+q)^2=2\xi P\cdot q-Q^2=Q^2(\xi-x)/x;
\end{equation}
领头阶 $\xi=x$，被击夸克沿 PDF Wilson 线的方向 $n_B=(0^+,1^-,\mathbf 0_T)$ 飞出。对固定 $x<1$ 的单个割图，领头钳夹面由靶共线子图、硬子图、沿 $n_B$ 的“流喷注”与软子图组成（\cref{fig:reduced-dis}）；
在非物理区的\emph{前向振幅}中后两者都不出现。

\paragraph{非物理区中的领头区域。}
对 $|\omega|<1$ 取参考系 $q^+=0$、$|\mathbf q_T|=Q$，于是 $2P\cdot q=2P^+q^-=\omega Q^2$，这一参数化对复 $\omega$ 也成立。
设 $k=\xi P+\kappa$ 是从靶共线子图流入硬部分的动量，$\kappa^+=0$，$|\boldsymbol\kappa_T|\lesssim\lambda Q$，$|\kappa^-|\lesssim\lambda^2Q^2/P^+$；由 $\kappa\cdot q=-\boldsymbol\kappa_T\cdot\mathbf q_T$，任何携带 $q$ 的线满足
\begin{equation}\label{eq:disall-unphysical-offshell}
(k+q)^2=Q^2(\xi\omega-1)+2\kappa\cdot q+k^2=Q^2(\xi\omega-1)+\order{\lambda Q^2},\qquad |\xi\omega-1|\ge1-|\omega|\quad(|\xi|\le1).
\end{equation}
所以对 $|\omega|\le\rho<1$，携带 $q$ 的线离壳至少 $Q^2(1-\rho)$，在每个区域中都是硬线，附着在它们上的软线随之进入硬子图（附着在硬子图上的软线被幂次压低，\cref{ch:ir}），
只附着在靶共线子图上的软线属于共线子图内部。领头区域因此只由靶共线子图 $J$ 与硬子图 $H$ 组成，二者由两条物理极化的部分子外加任意多条标量（纵向）极化的 K 胶子相连；
K 胶子对 $H$ 的附着之和经 Ward 恒等式成为 PDF 的 Wilson 线（一般情形是\cref{lem:method-ward-multi}的内容，全阶论证，[A]），其种子见\cref{fig:disall-seed}。

%%% BEGIN fig:disall-seed
\begin{figure}[htbp]
\centering
\begin{tikzpicture}[fd,
  fb lab/.style={font=\footnotesize,inner sep=1.5pt},
  fb op/.style={font=\large},
  fb hregion/.style={hard blob,fill=gray!22,draw=gray!60!black,line width=0.6pt,dash pattern=on 2.5pt off 1.5pt}]
  % ======================= left: one K gluon attached to the hard part =======================
  \begin{scope}
    \coordinate (V) at (-0.4,1.55);       % photon vertex
    \coordinate (G) at (1.0,2.3);         % attachment of the K gluon on the struck quark
    \coordinate (E) at (2.05,2.86);       % outgoing struck quark
    \begin{scope}[on background layer]
      \draw[fb hregion,rotate around={28.2:(0.3,1.92)}] (0.3,1.92) ellipse[x radius=1.12,y radius=0.42];
    \end{scope}
    \node[jet blob,ellipse,minimum width=1.9cm,minimum height=0.7cm,inner sep=0pt] (A) at (0,0) {};
    \node[font=\small] at (-0.45,0) {$A$};
    \draw[photon] (-1.75,2.3) -- (V);
    \draw[quark] (-0.4,0.34) -- (V);
    \draw[quark] (V) -- (G);
    \draw[quark] (G) -- (E);
    \draw[longit] (0.6,0.33) to[out=72,in=-108] (G);
    \fill[blue!60!black] ($(G)+(-0.06,-0.19)$) circle[radius=1.7pt];
    \node[vertex] at (V) {}; \node[vertex] at (G) {};
    \path[mom'={$q$},mom sep=0.24cm,mom length=0.45cm] (-1.75,2.3) -- (V);
    \path[mom={$l$},mom sep=0.24cm,mom length=0.45cm] (-0.4,0.34) -- (V);
    \path[mom={$p'-k$},mom sep=0.3cm,mom length=0.42cm,mom pos=0.42] (V) -- (G);
    \path[mom={$p'$},mom sep=0.24cm,mom length=0.42cm] (G) -- (E);
    \path[mom'={$k$},mom sep=0.26cm,mom length=0.45cm,mom pos=0.45] (0.6,0.33) to[out=72,in=-108] (G);
    \node[fb lab,text=gray!35!black] at ($(G)+(0.42,-0.36)$) {$H$};
    \node[fb lab,anchor=north] at (0.08,-0.42) {$k_\mu H^\mu_a=g\,t^a\,H_0$};
    \node[font=\large,anchor=east] at (-1.75,1.2) {$\displaystyle\sum_{H\text{ 内的附着}}$};
  \end{scope}
  \node[fb op] at (2.75,1.4) {$=$};
  % ======================= middle: H_0 =======================
  \begin{scope}[xshift=4.55cm]
    \coordinate (V) at (-0.25,1.65);
    \draw[photon] (-1.3,2.35) -- (V);
    \draw[quark] (-0.25,0.6) -- (V);
    \draw[quark] (V) -- (1.15,2.4);
    \node[vertex] at (V) {};
    \path[mom={$l+k$},mom sep=0.24cm,mom length=0.45cm] (-0.25,0.6) -- (V);
    \path[mom={$p'$},mom sep=0.24cm,mom length=0.42cm] (V) -- (1.15,2.4);
    \node[fb lab] at (0.6,1.42) {$H_0$};
  \end{scope}
  \node[fb op] at (6.15,1.4) {$\times$};
  % ======================= right: the Wilson-line vertex of the PDF operator =======================
  \begin{scope}[xshift=7.75cm]
    \coordinate (O) at (-0.35,1.65);     % operator vertex psi(0)
    \coordinate (W) at (1.75,2.8);       % towards infinity along n_B
    \coordinate (Y) at ($(O)!0.5!(W)$);
    \node[jet blob,ellipse,minimum width=1.9cm,minimum height=0.7cm,inner sep=0pt] (A) at (0,0) {};
    \node[font=\small] at (-0.45,0) {$A$};
    \draw[quark] (-0.35,0.34) -- (O);
    \draw[wilson=0.78] (O) -- (W);
    \node[opvertex] at (O) {};
    \draw[longit] (0.55,0.33) to[out=70,in=-110] (Y);
    \node[vertex] at (Y) {};
    \path[mom'={$k$},mom sep=0.26cm,mom length=0.45cm,mom pos=0.45] (0.55,0.33) to[out=70,in=-110] (Y);
    \node[fb lab,anchor=east] at ($(O)+(-0.12,0.05)$) {$\psi(0)$};
    \node[fb lab,anchor=west] at (W) {$\infty n_B$};
    \node[fb lab,anchor=north west] at ($(Y)+(0.12,-0.06)$) {$\ii g\,t^a n_B^\alpha$};
    \node[fb lab,anchor=south east] at ($($(O)!0.5!(Y)$)+(-0.02,0.12)$) {$\dfrac{-\ii}{n_B\cdot k-\ii0}$};
    \node[fb lab,anchor=north] at (0.4,-0.42) {$W^{(1)}[\infty n_B,0]$：未来指向};
  \end{scope}
  % summary formula
  \node[font=\footnotesize,anchor=north] at (4.4,-1.15) {$H_\mu\,\dfrac{k^\mu n_B^\alpha}{n_B\cdot k-\ii0}=\Big[\ii g\,t^an_B^\alpha\,\dfrac{-\ii}{n_B\cdot k-\ii0}\Big]H_0=W^{(1)}[\infty n_B,0]\,H_0$};
\end{tikzpicture}
\caption{DIS 中 PDF Wilson 线的种子（VF-WL-8，\cref{prop:method-ward-tree}）：一个来自靶共线子图 $A$ 的 K 胶子（蓝色，流入硬部分的动量为 $k$，圆点标出带标量极化 $k_\mu$ 的硬端）对硬部分 $H$ 的所有附着之和（领头阶即附着在末态被击夸克 $p'$ 上），
经 Ward 恒等式 $k_\mu H^\mu_a=g\,t^aH_0$ 化为去掉该胶子、入射动量为 $l+k$ 的硬部分 $H_0$，乘以夸克场 $\psi(0)$ 处出发、沿 $n_B$ 的未来指向 Wilson 线 $W^{(1)}[\infty n_B,0]$ 的单胶子项（双线箭头为路径方向 $0\to\infty n_B$）。
程函传播子中的 $-\ii0$ 继承自被替代的被击夸克传播子：$(p'-k)^2+\ii0\simeq-2p'^-k^++\ii0$ 的极点位于 $k^+$ 上半平面；Drell--Yan 中同样的附着落在入射反夸克上，给出 $+\ii0$。}
\label{fig:disall-seed}
\end{figure}
%%% END fig:disall-seed

把这一标度结构变成余项的精确陈述需要附录~N 的减除框架，回到物理区则需要区分两条途径。

\begin{proposition}[从非物理区前向振幅到矩与 $x$ 空间的两条途径 \tagR]\label{prop:disall-two-routes}
（$R_{\mathrm{cond}}$：在所列编号假设之下证明完整；结论的地位不高于其中最弱的假设。）前提：\cref{asm:appN-collins}（在复圆盘 $|\omega|\le\rho<1$ 上一致）、\cref{asm:appN-contour}、\cref{asm:appN-envelope}、\cref{asm:appN-separation}、\cref{asm:appN-cutsum}（用于 (c)）。（本命题只涉及部分子态。）
设部分子靶 $j$ 带有红外正规化（参数 $\lambda$）。在 $\as$ 的固定阶，对前向振幅 $T^{\mu\nu}(\omega)$ 的全部领头幂次图按\cref{def:method-subtraction}构造减除项 $C_RT^{\mu\nu}$，
记余项 $r^{\mu\nu}\equiv T^{\mu\nu}-\sum_RC_RT^{\mu\nu}$，$r_a$（$a=1,2$）为其不变振幅（以\cref{def:disall-breit}的投影提取）。
\begin{enumerate}[label=(\alph*),leftmargin=2em,itemsep=1pt,topsep=2pt]
\item 对 $|\omega|<1$，$T^{\mu\nu}$ 的领头区域只含靶共线模式与硬模式：没有末态，携带 $q$ 的线离壳 $\order{Q^2}$，只附着在靶共线子图上的软线属于共线子图。
因此只用到硬—共线近似子，交封闭的区域集合只有一个共线方向；并且对每个实数 $\omega\in[-\rho,\rho]$，$\rho<1$，有 $r^{\mu\nu}(\omega)=\order{\lambda^aL^K}$，
$a>0$ 与 $K$ 不依赖于 $\lambda$，常数对 $\varepsilon\in(0,\varepsilon_0]$ 一致。
\item 若 (a) 的界对复圆盘 $|\omega|\le\rho$ 一致成立，则每个矩都以 $\order{\rho^{-n}\lambda^aL^K}$ 的余项因子化：记 $M^{(2)}_n=\int_0^1\dd x\,x^{n-2}F_2$，
$M^{(1)}_n=\int_0^1\dd x\,x^{n-1}F_1$（$n$ 为偶数），$[\omega^k]f$ 为 $f$ 在 $\omega=0$ 处的第 $k$ 个 Taylor 系数，则
\begin{equation}\label{eq:disall-cauchy}
\begin{gathered}
\Big|M^{(2)}_n-\tfrac14[\omega^{n-1}]\textstyle\sum_RC_RT_2\Big|\le\tfrac14\rho^{1-n}\max_{|\omega|=\rho}|r_2(\omega)|,\\
\Big|M^{(1)}_n-\tfrac14[\omega^{n}]\textstyle\sum_RC_RT_1\Big|\le\tfrac14\rho^{-n}\max_{|\omega|=\rho}|r_1(\omega)|.
\end{gathered}
\end{equation}
仅在实线段上的界不给出 Taylor 系数的界：复圆盘上的一致性是额外的输入。
\item 实线段上的界不控制不连续部分。固定 $x<1$ 处的 $x$ 空间陈述需要在\cref{asm:appN-cutsum}之下对割图求和后的被积函数作减除分析；当 $x\to1$ 时这一分析中的常数不一致（阈值区）。
\end{enumerate}
\uses{def:disall-breit,def:dis-compton,def:method-subtraction,def:method-collinear-approx,def:method-collins-conditions,asm:appN-collins,asm:appN-cutsum,asm:appN-contour,asm:appN-envelope,asm:appN-separation,thm:appN-subtraction,lem:appN-error-table,lem:method-local-accuracy,lem:method-meets,prop:dis-compton-analytic,prop:dis-dispersion}
% errtype: E-N
\end{proposition}

\begin{proof}
(a) 由\cref{eq:disall-unphysical-offshell}，携带 $q$ 的线都离壳 $\ge Q^2(1-\rho)-\order{\lambda Q^2}$，所以每个领头区域是在每条线上指派 $A$（靶共线）或 $H$，携带 $q$ 的线必为 $H$。
区域按逐线偏序比较，交逐线取得（\cref{lem:method-meets}）：钳夹集相交时 $\pi_{R\wedge R'}=\pi_R\cap\pi_{R'}$，交属于交封闭的区域集合；不相交时二者没有公共下界，(CS2) 须直接验证。
$T_R$ 在 $H_R$ 中作 $k\to\pi_+(k)=(k^+,0,\mathbf 0_T)$ 并在界面插入 Dirac 投影 $\frac12\gamma^-\gamma^+$ 或胶子投影（\cref{def:method-collinear-approx}，只用误差表中共线近似的各行）；它不引入新分母，
被投影的硬分母在 $U_R$ 上由\cref{asm:appN-separation}(i) 满足 $|l^2-M_l^2|\ge c\,t_R^2$，非物理区又没有末态割线，因此 $C_RT^{\mu\nu}$ 在\cref{asm:appN-contour}的围道上没有新奇点，不出现软近似分母的 $\ii0$ 相容性问题。
极小区域的 (CS1) 由\cref{lem:method-local-accuracy}给出（前提\cref{asm:appN-envelope,asm:appN-separation}，误差各行见\cref{lem:appN-error-table}），其余的 (CS1)(CS2)(CS3) 是\cref{asm:appN-collins}的内容，
于是\cref{thm:appN-subtraction}（前提\cref{asm:appN-contour}(a)）给出 $\int|r|\le C\lambda^aL^K$，对 $\varepsilon\in(0,\varepsilon_0]$ 一致；\cref{asm:appN-collins}在圆盘上一致，故常数对 $\omega\in[-\rho,\rho]$ 一致。

(b) 在固定阶，$T_{1,2}(\omega)$ 在去掉实轴上 $|\omega|\ge1$ 两条割线的平面上解析（\cref{prop:dis-compton-analytic}）；$C_RT^{\mu\nu}$ 只通过入射动量为 $\xi p$（$|\xi|\le1$）的被投影硬因子依赖于 $\omega$，
其阈值在 $(\xi p\pm q)^2=0$ 即 $\xi\omega=\pm1$，故 $r_a$ 在 $|\omega|<1$ 中解析，Cauchy 估计给出 $|[\omega^k]r_a|\le\rho^{-k}\max_{|\omega|=\rho}|r_a|$。由色散关系（\cref{prop:dis-dispersion}）
$T_2=\sum_{n\ \text{偶}}4\omega^{n-1}M^{(2)}_n$，$T_1-T_1(0)=\sum_{n\ \text{偶}}4\omega^nM^{(1)}_n$，即得\cref{eq:disall-cauchy}。反之，$\lambda^a\sin(\omega/\lambda)$ 在实轴上不超过 $\lambda^a$，
其奇数阶 Taylor 系数却为 $\pm\lambda^{a-k}/k!$，对 $k>a$ 随 $\lambda\to0$ 发散。

(c) 不连续部分的逐点大小不受 $|\omega|<1$ 上取值的控制：若某余项的不连续部分为 $\sin(x/\lambda^a)$（$0\le x\le1$），分部积分给出 $|\int_0^1\dd x\,x^n\sin(x/\lambda^a)|\le2\lambda^a$（一切 $n\ge0$），
由色散关系它在 $|\omega|\le\rho$ 上对应的前向振幅不超过 $8\rho\lambda^a/(1-\rho^2)$，而它本身逐点为 $\order{1}$（同类反例见\cref{rem:dis-nonuniform}）。所以固定 $x$ 须直接处理割图：
在\cref{asm:appN-cutsum}之下，每个前向图 $G$ 的割图之和在公共区域 $\Omega_G$ 上有公共被积函数 $F_G$（光锥序表示），近似子与割图求和对易，(GA1)–(GA4)、(CS1)–(CS3) 对 $F_G$ 成立，
\cref{thm:appN-subtraction}便可在固定 $x$ 处应用于 $F_G$；其中不变质量为 $Q^2(\xi-x)/x$ 的末态整体并入硬子图，内部的末态共线与软钳夹在割图求和后相消。
当 $\xi\to x$ 时这一不变质量趋于零，流喷注成为真正的共线子图；对 $1-x\lesssim\lambda^b$（某个 $b>0$）不存在使流喷注为硬的区域，常数随 $x\to1$ 增长：这一端点正是加号分布与阈值区的来源（\cref{ch:thr}）。
\end{proof}

\noindent
(b) 是“矩的途径”：不需要 OPE，但需要复 $\omega$ 上的一致界，只给出分布意义下的结论；(c) 是“$x$ 空间的途径”，代价是割图求和的公共表示\cref{asm:appN-cutsum}，主定理\cref{thm:disall-factorization}走这一条路。

% ---------------------------------------------------------------------
\section{末态相互作用与幺正性}\label{sec:disall-fsi}

在单个割图中，被击夸克与靶残余之间可以交换任意多个胶子。它们不破坏因子化的原因是幺正性：结构函数只依赖于对全部末态求和的结果，即前向振幅的不连续部分，而前向振幅中携带光子动量的线不被钳夹。

\begin{lemma}[末态喷注的非钳夹性 \tagA]\label{lem:disall-fsi}
对末态的全部割图求和给出前向振幅的不连续部分，$\Imag T_a(\omega+\ii0)=2\pi F_a$（$\omega>1$）。在前向振幅中，携带大分量 $q^-$ 的线不被钳夹：
在靶共线钳夹面附近，经过这些线的圈动量的负分量围道可以离开它们的极点一段 $\order{Q^2/P^+}$ 的距离，所以流喷注以及附着在它上面的软线都属于硬子图 $H$。
由\cref{eq:disall-current-virtuality}，$(\xi P+q)^2=Q^2(\xi-x)/x$：对 $\xi-x=\order{1}$，末态系统的不变质量为 $\order{Q^2}$，流喷注并入 $H$；
只在卷积的端点 $\xi\to x$ 处它才成为真正的共线喷注，这一端点产生系数函数中的加号分布，也是阈值区的来源。
\uses{prop:method-fsi,def:disall-breit,prop:disall-two-routes,prop:dis-compton-analytic}
% errtype: E-S
\end{lemma}

\begin{proofsketch}
(1) 一个前向图的全部末态割线之和等于它的不连续部分；按\cref{ch:dis}的归一化，这就是 $\Imag T_a=2\pi F_a$（\cref{prop:dis-compton-analytic}）。
(2) 携带 $q$ 的线 $k+q$ 的分母作为 $k^-$ 的函数，极点在 $k^-\simeq-q^-=-Q^2/(2x_NP^+)$ 附近，与共线钳夹面 $k^-=\order{\lambda^2Q^2/P^+}$ 相距 $\order{Q^2/P^+}$；
非物理区中的定量形式即\cref{eq:disall-unphysical-offshell}（\cref{prop:disall-two-routes}(a)）。附着在流喷注上的软动量 $\ell$ 只通过 $\ell^+$ 耦合（$(p'+\ell)^2-p'^2\simeq2p'^-\ell^+$），
极点都在 $\ell^+$ 平面的同一侧，围道可以移开：这是\cref{prop:method-fsi}的 DIS 情形。(3) 当 $\xi\to x$ 时流喷注的虚度趋于零，额外的末态共线区域给出加号分布（\cref{sec:disall-nlo}），
在固定 $x<1$ 处对光滑的 PDF 可积。完整论证见\cite[第~8.7~节]{CollinsSoperSterman1989}与\cite{LibbySterman1978a}。
\end{proofsketch}

\paragraph{地位与缺口}
这是 E-S 型论证：它说明末态钳夹面不是领头的，但不给出余项的积分界。固定 $x$ 处所需的割图版本建立在\cref{asm:appN-cutsum}之上（\cref{prop:disall-two-routes}(c)）；
一般的相消机制\cref{prop:method-fsi}（全阶论证，[A]）还依赖光锥序表示的假设\cref{asm:appO-lcopt}。文献中这一步是标度陈述：\cite[第~8.7~节]{CollinsSoperSterman1989}写道
「“the sum over final states eliminates pinch singularities involving final state jets”」，即对末态求和消除了涉及末态喷注的钳夹奇点。对 DIS 的割图证明\cref{asm:appN-cutsum}即可把本引理提升为 $R_{\mathrm{cond}}$ 意义下的陈述。

\begin{proposition}[DIS 的领头区域 \tagA]\label{prop:disall-regions}
$T^{\mu\nu}$ 的领头区域（等价地，对割图求和后 $W^{\mu\nu}$ 的领头区域）由一个靶共线子图与一个硬子图组成：二者由两条部分子（两条夸克线，或两条横向极化的胶子）外加任意多条 $A^+$ 极化（K）胶子相连；
没有软子图；对割图求和后没有末态共线子图。（本命题只涉及部分子态；强子版本即\cref{asm:intro-bridge}(Br1)。）
\uses{thm:ir-leading-regions,lem:disall-fsi}
% errtype: E-S
\end{proposition}

\begin{proofsketch}
(1) 对每个割图应用\cref{thm:ir-leading-regions}（全阶论证，[A]）：领头区域由靶共线子图 $J_A$、流喷注 $J_q$ 与只以胶子附着在二者上的软子图组成，没有软线附着在硬子图上（\cref{fig:reduced-dis}）。
(2) 由\cref{lem:disall-fsi}，割图求和后 $J_q$ 与附着在其上的软线并入硬子图；附着在硬子图上的每条软胶子压低 $\lambda^2$，只连接 $J_A$ 自身的软线是共线子图内部的子区域。
(3) 活动部分子为胶子时，横向胶子与 K 胶子组合为 $k^+A^j-k^jA^+$，即胶子 PDF 算符中 $G^{+j}$ 的 Fourier 分量。文献：\cite{LibbySterman1978a}，\cite[第~8.7~节]{CollinsSoperSterman1989}。
\end{proofsketch}

\paragraph{地位与缺口}
这是标度陈述（E-S）：嵌套区域没有 Minkowski 空间中与 Weinberg 定理相当的结果，个别 Feynman 规范图中的超领头项只在对图求和后相消。
其积分形式是\cref{asm:appN-envelope}（幂次计数包络）与\cref{asm:appN-separation}（分离性），作为\cref{thm:disall-factorization}的前提出现；固定 $x$ 处的割图版本依赖于\cref{asm:appN-cutsum}。

% ---------------------------------------------------------------------
\section{减除项的因子化结构、重整化相容性与两级梯形图}\label{sec:disall-subtraction}

减除定理（\cref{thm:method-subtraction}，全阶论证，[A]）给出 $\Gamma=\sum_RC_R\Gamma+r$，$r\in\mathfrak S(\Omega)$，但它本身不说明 $\sum_RC_R\Gamma$ 具有“硬部分 $\otimes$ PDF”的形式。
DIS 只有一个共线方向、没有软区域，较小区域的共线子图\emph{包含}较大区域的共线子图；这使减除项具有嵌套的乘积结构，并可以与 $\MSbar$ 重整化逐阶对接。

\begin{proposition}[DIS 减除项的因子化结构 \tagR]\label{prop:disall-factorized-subtraction}
（$R_{\mathrm{cond}}$：在所列编号假设之下证明完整；结论的地位不高于其中最弱的假设。）前提：\cref{asm:appN-ward-product}、\cref{asm:appN-uv}（及\cref{prop:disall-regions}）。（本命题只涉及部分子态。）
设 $\Gamma$ 遍历部分子靶 $j$ 的前向振幅在 $\as$ 某一固定阶 $n$ 的全部领头幂次图，区域如\cref{prop:disall-regions}，动量分配是适配的（每个区域的界面动量与其硬子图内部的圈动量彼此分开）。则
\begin{enumerate}[label=(\roman*),leftmargin=2em,itemsep=1pt,topsep=2pt]
\item 若 $R'<R$，则 $J_{R'}\supset J_R$，$H_{R'}\subset H_R$；因此 $T_{R'}$ 只改变 $H_R$ 的因子以及 $J_R$ 与 $H_R$ 之间的投影，从不改变 $J_R$ 内部的因子；
\item $C_R\Gamma=T_R\prod_{R'<R}(1-T_{R'})\Gamma$，即“减去全部较小区域之后的硬因子在投影动量处的值”与“精确的 $J_R$”之积；
\item 对全部图与全部区域求和，并按共线部分的阶 $n_J$ 分组，
\begin{equation}\label{eq:disall-factorized}
\sum_\Gamma\ \sum_{R:\ \mathrm{ord}(J_R)=n_J}C_R\Gamma=\sum_i\int\frac{\dd\xi}{\xi}\,\hat{\mathsf H}^{[n-n_J]}_i(\xi p,q)\,f^{\mathrm{bare},[n_J]}_{i/j}(\xi),
\end{equation}
其中 $f^{\mathrm{bare}}_{i/j}$ 是\cref{def:pdf-quark,def:pdf-gluon}的算符在部分子态 $j$ 中的 $D$ 维裸矩阵元（不作算符重整化），$\hat{\mathsf H}_i$ 只依赖于投影后的部分子动量 $\xi p$，
与靶无关；色投影 $\Trc[HJ]\simeq\frac1{N_c}\Trc H\,\Trc J$（胶子为 $\frac1{N_c^2-1}$）。特别地，$T_{(j)}-r_{(j)}=\sum_i\hat{\mathsf H}^{\mathrm{bare}}_i\otimes f^{\mathrm{bare}}_{i/j}$。
\end{enumerate}
\uses{prop:method-product,def:method-subtraction,prop:disall-regions,def:pdf-quark,def:pdf-gluon,asm:appN-ward-product,asm:appN-uv,lem:appN-algebra}
% errtype: E-W
\end{proposition}

\begin{proof}
(i) 由\cref{prop:disall-regions}，每个领头区域由包含两个光子顶点的硬子图 $H_R$ 确定，$J_R$ 是其余部分；逐线偏序中 $A<H$，故 $R'\le R$ 当且仅当 $H_{R'}\subseteq H_R$，即 $J_{R'}\supseteq J_R$。
$T_{R'}$ 在 $H_{R'}$ 中把界面动量换成 $\pi_+(k)$（适配的动量分配下内部圈动量不变），只改变 $H_{R'}\subset H_R$ 的因子；它在 $R'$ 的界面线上插入的投影，若该线在 $R$ 中为硬则位于 $H_R$ 内部，
若该线属于 $J_R$ 则它也是 $R$ 的界面线，投影与 $T_R$ 的相同（$\frac12\gamma^-\gamma^+$ 或 $\mathcal P_A$），由幂等性合而为一。所以 $J_R$ 内部的因子不变。

(ii) 只有一个共线方向，没有方向翻转的线；适配的动量分配下各近似子是坐标投影（动量分量置零加投影），幂等且满足 (Com)(Abs)（\cref{lem:appN-algebra}），区域集合取为交封闭的（附录~N 的约定 (GA2)），
于是乘积公式\cref{prop:method-product}给出 $C_R\Gamma=T_R\prod_{R'<R}(1-T_{R'})\Gamma$。写 $\Gamma=H_R\otimes J_R$（对界面指标缩并、对界面动量积分），由 (i) 得 $\prod_{R'<R}(1-T_{R'})\Gamma=\tilde H_R\otimes J_R$，
$\tilde H_R\equiv\prod_{R'<R}(1-T_{R'})H_R$，故 $C_R\Gamma=\int\frac{\dd^Dk}{(2\pi)^D}\Tr[\tilde H_R(\pi_+k)\frac14\gamma^-]\Tr[\gamma^+J_R(k)]$（胶子界面以 $\mathcal P_A$ 代替）。

(iii) $\tilde H_R(\pi_+k)$ 只依赖于 $k^+=\xi p^+$，$k^-$、$\mathbf k_T$ 的积分只作用于 $J_R$ 并延伸到全部取值。以 $J_R(k)=\int\dd^Dw\,\ee^{\ii k\cdot w}\langle j|\bar\psi(0)\psi(w)|j\rangle$
（含 K 胶子与 Wilson 线）计，由 $\int\dd k^-\dd^{D-2}\mathbf k_T\,\ee^{\ii k\cdot w}=(2\pi)^{D-1}\delta(w^+)\delta^{D-2}(\mathbf w_T)$、平移不变性与\cref{def:pdf-quark}的归一化 $\int\dd w^-/(4\pi)$，
\begin{equation}\label{eq:disall-chain}
\int\frac{\dd k^-\,\dd^{D-2}\mathbf k_T}{(2\pi)^D}\,\Tr\big[\gamma^+J_R(k)\big]=\frac1{2\pi}\int\dd w^-\,\ee^{-\ii\xi p^+w^-}\langle j|\bar\psi(w)\gamma^+W[w,0]\psi(0)|j\rangle=2f_{q/j}(\xi);
\end{equation}
又 $\dd^Dk=p^+\dd\xi\,\dd k^-\dd^{D-2}\mathbf k_T$，故 $C_R\Gamma=\int\frac{\dd\xi}{\xi}\frac12\Tr[\tilde H_R(\xi p)\xi\sla p]f(\xi)$，其中 $\xi\sla p=\sum_su\bar u$，$\frac12\Tr[\cdots]$ 是对在壳部分子 $\xi p$ 自旋平均的硬振幅。
对全部图求和时，K 胶子对 $\tilde H_R$ 的附着由\cref{asm:appN-ward-product}(WI1) 分离为沿 $n_B$ 的 Wilson 线（(WI1) 明确包括已被其他近似子改变过的被积函数），与 $J_R$ 合成\cref{def:pdf-quark,def:pdf-gluon}中算符的矩阵元，
(WI2) 保证按共线部分的阶分组后得到\cref{eq:disall-factorized}；\cref{asm:appN-uv}(U1) 使每个图在被积函数层次上已作紫外重整化，(U2) 保证 $T_R$ 与紫外减除只差不领头的多项式。
色平均的靶有 $\langle J_{ab}\rangle\propto\delta_{ab}$，K 胶子并入 Wilson 线后算符是色单态，故色投影对整个乘积成立。$\hat{\mathsf H}^{[n_H]}_i$ 定义为全部硬拓扑上 $\frac12\Tr[\tilde H_R(\xi p)\xi\sla p]$（胶子对 $D-2$ 个极化平均）之和，只依赖于部分子种类与投影动量，因而与靶无关。
\end{proof}

\noindent
若较小区域使 $J_R$ 中的线变软（事件形状、Drell--Yan、TMD），$T_{R'}$ 会改变 $J_R$ 的因子，裸共线因子须换成“共线因子减去其零箱”（单圈离壳三角图中 $C_{C_1}\Gamma=T_{C_1}\Gamma-T_{C_1}T_S\Gamma$）；
DIS 没有领头的软区域，这正是“DIS 是最简单情形”的精确含义。下面接上 $\MSbar$ 重整化。按\cref{asm:appN-multren}(Z1) 写 $f^{\mathrm{bare}}_N=Z_Nf^{\mathrm{ren}}_N$（部分子种类空间中的矩阵）；它与\cref{def:pdf-quark}中 $f_{\mathrm{ren}}=Z\otimes f_{\mathrm{bare}}$ 的 $Z$ 互为逆，
单圈 $Z_N=1+\frac{\as}{2\pi}\frac1\varepsilon\gamma(N)+\order{\as^2}$，$\gamma_{ij}(N)=\int_0^1\dd z\,z^{N-1}P_{ij}$。

\begin{theorem}[减除与 $\overline{\rm MS}$ 重整化的相容性（矩空间） \tagR]\label{thm:disall-renorm-compat}
（$R_{\mathrm{cond}}$：在所列编号假设之下证明完整；结论的地位不高于其中最弱的假设。）前提：\cref{asm:appN-multren}（(Z1)、(Z3)(i)）、\cref{asm:appN-ir-regulator}、\cref{asm:appN-contour}（(b)）（及\cref{prop:disall-two-routes}、\cref{prop:disall-factorized-subtraction}的前提）。（本定理只涉及部分子态。）
本定理的内容等级为 [C]（经由\cref{asm:appN-ir-regulator}）：它是正规子路线下的严格一致性陈述，任何全阶论证（[A]）的陈述都不以它为依赖；DIS 与其他过程的定理使用维数正规化路线\cref{asm:appN-multren}(Z3)(ii)。
考虑非物理区前向振幅的第 $N$ 阶矩。由\cref{prop:disall-two-routes}(a)(b)与\cref{prop:disall-factorized-subtraction}，
\begin{equation}\label{eq:disall-renorm-start}
T_{(j),N}=\sum_i\hat{\mathsf H}^{\mathrm{bare}}_{i,N}\,f^{\mathrm{bare}}_{i/j,N}+r_{(j),N},\qquad |r_{(j),N}|\le C\lambda^aL^K .
\end{equation}
\begin{enumerate}[label=(\alph*),leftmargin=2em,itemsep=1pt,topsep=2pt]
\item （前提 (Z1)、\cref{asm:appN-ir-regulator}）$C_{i,N}\equiv\sum_k\hat{\mathsf H}^{\mathrm{bare}}_{k,N}Z_{N,ki}$ 在 $\varepsilon=0$ 处有限；
\item $C_{i,N}$ 与靶及其红外正规化无关；（前提 (Z1)、(Z3)(i)）它等于维数正规化中部分子匹配系数的矩（\cref{prop:pdf-partonic-dr}的构造）；
\item 逐阶地 $C^{[n]}=\lim_{\lambda\to0}\big[T^{[n]}_{(j)}-\sum_{m=0}^{n-1}C^{[m]}f^{\mathrm{ren}[n-m]}_{\cdot/j}\big]$：每一项都像 $\ln\lambda$ 一样发散，只有组合有极限；以维数正规化作红外正规化时 $\lambda\to0$ 一步可以略去。
\end{enumerate}
\uses{prop:disall-factorized-subtraction,lem:method-no-pole,prop:pdf-partonic-dr,def:prelim-msbar,asm:appN-multren,asm:appN-ir-regulator,asm:appN-contour,prop:disall-two-routes,cor:method-removing-eps}
% errtype: E-N
\end{theorem}

\begin{proof}
(a) 由 (Z1)，$\hat{\mathsf H}^{\mathrm{bare}}f^{\mathrm{bare}}_{\cdot/j}=\hat{\mathsf H}^{\mathrm{bare}}Zf^{\mathrm{ren}}_{\cdot/j}=Cf^{\mathrm{ren}}_{\cdot/j}$（矩阵记号），代入\cref{eq:disall-renorm-start}得
$C=\hat{\mathsf H}^{\mathrm{bare}}Z=(T_{(j)}-r_{(j)})(f^{\mathrm{ren}}_{\cdot/j})^{-1}$。左边与 $\lambda$ 及靶无关（$\hat{\mathsf H}^{\mathrm{bare}}$ 只依赖于投影动量，$Z$ 是与靶无关的纯 $\MSbar$ 极点级数），
且在 $\varepsilon=0$ 附近亚纯（每个 $\int C_R\Gamma$ 由\cref{asm:appN-contour}(b) 亚纯延拓）。右边在 $(0,\varepsilon_0]$ 上有界：$T_{(j)}$ 是守恒流关联函数的矩，紫外有限，红外由 $\lambda$ 正规化；
$r_{(j)}=T_{(j)}-\sum_R\int C_R\Gamma$ 中后者由\cref{cor:method-removing-eps}在 $\varepsilon=0$ 处没有极点，且 $|r_{(j)}|\le C\lambda^aL^K$ 对 $\varepsilon$ 一致；
由\cref{asm:appN-ir-regulator}，$f^{\mathrm{ren}}_{N,\cdot/j}$ 在 $\varepsilon=0$ 处有限并逐阶可逆，且该正规化保持\cref{prop:disall-factorized-subtraction}所用的 Ward 恒等式并使\cref{prop:disall-two-routes}(a)(b) 成立。
在 $\varepsilon=0$ 附近亚纯、在 $(0,\varepsilon_0]$ 上有界的函数没有极点（\cref{lem:method-no-pole}），故 $C$ 在 $\varepsilon=0$ 处有限。

(b) $C$ 由左边定义，因此与靶及 $\lambda$ 无关。对在壳无质量部分子靶、以维数正规化为唯一的红外正规化，(Z3)(i) 断言因子化恒等式以同一个 $\hat{\mathsf H}^{\mathrm{bare}}$ 成立，且 $f^{\mathrm{bare}}_{\cdot/j}=\delta$（无标度），
于是 $T^{\mathrm{DR}}_{(j)}=\hat{\mathsf H}^{\mathrm{bare}}$，$f^{\mathrm{ren,DR}}_{\cdot/j}=Z^{-1}$，$C=T^{\mathrm{DR}}_{(j)}(f^{\mathrm{ren,DR}}_{\cdot/j})^{-1}$：维数正规化的部分子结构函数乘以 $\MSbar$ 部分子 PDF 的逆，
即\cref{prop:pdf-partonic-dr}的构造（$\MSbar$ 的归一化见\cref{def:prelim-msbar}）。

(c) 把 $f^{\mathrm{ren}}_{\cdot/j}=\delta+\sum_{k\ge1}f^{\mathrm{ren}[k]}_{\cdot/j}$、$T_{(j)}=\sum_nT^{[n]}_{(j)}$、$C=\sum_nC^{[n]}$ 代入 $T_{(j)}=Cf^{\mathrm{ren}}_{\cdot/j}+r_{(j)}$，比较 $\as^n$ 的系数：
$C^{[n]}=T^{[n]}_{(j)}-\sum_{m<n}C^{[m]}f^{\mathrm{ren}[n-m]}_{\cdot/j}-r^{[n]}_{(j)}$。\cref{prop:disall-two-routes}(a) 的界在每一固定阶成立，$r^{[n]}_{(j)}\to0$（$\lambda\to0$），而左边不依赖于 $\lambda$，即得 (c)。
\end{proof}

\noindent
回到 $x$ 空间要用\cref{thm:pdf-dis-xspace}与附录~K 的唯一性。内容等级 [C] 的来源是：离壳外部部分子与 (Z1) 的在壳性矛盾，胶子质量破坏非阿贝尔 Ward 恒等式，在壳有质量夸克只正规化夸克靶而矩阵可逆性需要胶子靶，
迄今没有保持 Ward 恒等式又能正规化在壳胶子共线发散的红外正规化。过程定理因此改用 (Z3)(ii)，它在 NNLO 与 N$^3$LO 系数函数的有限性中被逐阶验证。
两级梯形图是上述结构的最小例子，也说明为什么只能对分组的表达式作估计（\cref{fig:disall-ladder-regions}）。

\begin{example}[两级梯形图：三个区域与分组余项]\label{ex:disall-ladder}
外夸克离壳（$p^2=-m^2$）、内线无质量的 DIS 两级梯形图：$\Gamma=\mathsf H_0(k_2)\,K_2(k_2,k_1)\,K_1(k_1,p)$，$\mathsf H_0$ 是含两个光子顶点的顶端夸克线，$K_2$、$K_1$ 是上、下两级横档（含脊线传播子），
$(k_1,k_2)$ 为两级横档之后的脊线动量。领头区域有\emph{三个}，构成链 $R_2<R_1<R_H$：$R_2$ 中两级横档都与 $p$ 共线，$R_1$ 中只有下级横档共线，$R_H$ 中一切都是硬的（靶被投影为无质量在壳部分子）。
$T_2$ 在 $\mathsf H_0$ 中作 $k_2\to\pi_+k_2$ 并插入 Dirac 投影 $\mathcal P$；$T_1$ 作 $k_1\to\pi_+k_1$ 并插入 $\mathcal P$（只改变 $K_2$）；$T_H$ 作 $p\to\hat p=\pi_+(p)$、$m\to0$（只改变 $K_1$）；它们彼此对易（\cref{prop:appN-ladder}）。
记 $\Delta\mathsf H_0\equiv\mathsf H_0(k_2)-\mathsf H_0(\pi_+k_2)\mathcal P$，$\Delta K_2\equiv K_2(k_2,k_1)-K_2(k_2,\pi_+k_1)\mathcal P$，$\Delta K_1\equiv K_1(k_1,p)-K_1(k_1,\hat p)$。
全序区域链的链展开（\cref{prop:method-chains}）给出 $r=(1-T_H)(1-T_1)(1-T_2)\Gamma$ 与
\begin{equation}\label{eq:disall-ladder}
\begin{gathered}
C_2\Gamma=[\mathsf H_0(\pi_+k_2)\mathcal P]\,K_2K_1,\qquad C_1\Gamma=\Delta\mathsf H_0\,[K_2(k_2,\pi_+k_1)\mathcal P]\,K_1,\\
C_H\Gamma=\Delta\mathsf H_0\,\Delta K_2\,K_1(k_1,\hat p),\qquad r=\Delta\mathsf H_0\,\Delta K_2\,\Delta K_1 .
\end{gathered}
\end{equation}
$C_2\Gamma$ 是 $\mathsf H^{[0]}\otimes f^{[2]}$；$C_1\Gamma$ 对 $k_2$ 积分得到单圈硬系数 $\hat{\mathsf H}^{[1]}(k_1^+)$（入射部分子为 $\pi_+k_1$ 的未减除单圈硬部分减去树图 $\otimes$ 单圈部分子共线因子），即 $\hat{\mathsf H}^{[1]}\otimes f^{[1]}$；
把 $\Delta\mathsf H_0\Delta K_2$ 展开为四项，依次是无质量靶的两圈梯形图 $T^{[2]}$、$\mathsf H^{[0]}\otimes f^{[2]}$、$\mathsf H^{[1],\mathrm{unsub}}\otimes f^{[1]}$ 与 $\mathsf H^{[0]}\otimes f^{[1]}\otimes f^{[1]}$，
合起来 $\int C_H\Gamma=T^{[2]}-\hat{\mathsf H}^{[1]}\otimes f^{[1]}-\mathsf H^{[0]}\otimes f^{[2]}$，即\cref{thm:disall-renorm-compat}(c) 在 $n=2$ 时的被积函数形式。
全硬区域 $R_H$ 携带两圈系数函数：只计入 $R_1$、$R_2$ 时，$\Gamma-C_2\Gamma-C_1\Gamma=C_H\Gamma+r$ 中的 $C_H\Gamma$ 是领头的。
分组的必要性：单个链项 $T_H(1-T_2)\Gamma=\Delta\mathsf H_0\,K_2(k_2,k_1)\,K_1(k_1,\hat p)$ 不属于 $L^1(U_{R_2})$——$k_1\parallel\hat p$ 时无质量的 $K_1(k_1,\hat p)$ 有 $\int\dd^{2-2\varepsilon}\mathbf k_{1T}/\mathbf k_{1T}^2$ 型的共线奇点
（只对 $\varepsilon<0$ 收敛），而 $\Delta\mathsf H_0K_2$ 在那里不为零；分组后的 $C_H\Gamma$ 多出的因子 $\Delta K_2$ 在 $k_1\to\pi_+k_1$ 时至少线性地趋于零，积分收敛。
所以误差估计只对分组的表达式成立，“减除定理可以逐项归纳证明”是错误的。上述代数是精确的（\cref{prop:appN-ladder}）；这一费米子梯形图的积分范数条件 (CS1)–(CS3) 没有被证明，标量模型中单梯级梯形图的完整证明见\cref{sec:appN-examples}。
\uses{thm:method-subtraction,prop:method-chains,def:method-collins-conditions,thm:disall-renorm-compat,prop:appN-ladder}
\end{example}

%%% BEGIN fig:disall-ladder-regions
\begin{figure}[htbp]
\centering
\begin{tikzpicture}[fd,
  fb lab/.style={font=\footnotesize,inner sep=1.5pt},
  fb kern/.style={font=\scriptsize,inner sep=1pt,text=gray!25!black},
  fb R2/.style={region,draw=blue!60!black,line width=0.7pt},
  fb R1/.style={region,draw=orange!75!black,line width=0.7pt},
  fb RH/.style={region,draw=gray!55!black,line width=0.7pt}]
  \def\fbX{1.2}
  % ======================= the two-rung ladder (cut graph; amplitude on the left) =======================
  % left rail (amplitude): p -> rung K_1 -> k_1 -> rung K_2 -> k_2 -> photon vertex
  \draw[quark] (-\fbX,0) -- (-\fbX,1.1);
  \draw[quark] (-\fbX,1.1) -- (-\fbX,2.3);
  \draw[quark] (-\fbX,2.3) -- (-\fbX,3.5);
  % right rail (conjugate amplitude)
  \draw[quark] (\fbX,3.5) -- (\fbX,2.3);
  \draw[quark] (\fbX,2.3) -- (\fbX,1.1);
  \draw[quark] (\fbX,1.1) -- (\fbX,0);
  % struck quark and photons (H_0)
  \draw[quark=0.25] (-\fbX,3.5) -- (\fbX,3.5);
  \draw[photon] (-1.62,4.72) -- (-\fbX,3.5);
  \draw[photon] (\fbX,3.5) -- (1.62,4.72);
  % rungs: real gluons crossing the cut
  \foreach \y in {1.1,2.3} {
    \draw[gluon] (-\fbX,\y) -- (\fbX,\y);
    \node[vertex] at (-\fbX,\y) {}; \node[vertex] at (\fbX,\y) {};}
  \node[vertex] at (-\fbX,3.5) {}; \node[vertex] at (\fbX,3.5) {};
  \draw[cut] (0,-0.35) -- (0,4.75);
  % momenta and kernels
  \path[mom'={$k_1$},mom sep=0.25cm,mom length=0.45cm] (-\fbX,1.1) -- (-\fbX,2.3);
  \path[mom'={$k_2$},mom sep=0.25cm,mom length=0.45cm,mom pos=0.42] (-\fbX,2.3) -- (-\fbX,3.5);
  \node[fb lab,anchor=north] at (-\fbX,-0.02) {$p$};
  \node[fb lab,anchor=north] at (\fbX,-0.02) {$p$};
  \node[fb lab,anchor=north,text=gray!30!black] at (0,-0.42) {$p^2=-m^2$};
  \node[fb lab,anchor=south east] at (-1.62,4.72) {$q$};
  \node[fb kern,anchor=south] at (0.6,1.25) {$K_1$};
  \node[fb kern,anchor=south] at (0.6,2.45) {$K_2$};
  \node[fb kern,anchor=south] at (-0.6,3.53) {$\mathsf H_0$};
  % nested hard subgraphs of the three regions
  \draw[fb R2] (-1.55,3.24) rectangle (1.55,3.98);
  \draw[fb R1] (-1.78,2.04) rectangle (1.78,4.14);
  \draw[fb RH] (-2.01,0.84) rectangle (2.01,4.3);
  \node[font=\scriptsize,text=blue!55!black,anchor=south,inner sep=1pt] at (0.72,3.25) {$R_2$};
  \node[font=\scriptsize,text=orange!65!black,anchor=south east,inner sep=1pt] at (1.76,2.06) {$R_1$};
  \node[font=\scriptsize,text=gray!35!black,anchor=south east,inner sep=1pt] at (1.99,0.86) {$R_H$};
  % ======================= the regions and the grouped remainder =======================
  \node[anchor=north west,align=left,font=\footnotesize] (T) at (2.65,4.3) {%
    区域链 $R_2<R_1<R_H$（虚线框为硬子图）：\\[3pt]
    {\color{blue!55!black}$R_2$}：两级横档都与 $p$ 共线；\\[1pt]
    {\color{orange!65!black}$R_1$}：上级横档为硬，下级横档共线；\\[1pt]
    {\color{gray!35!black}$R_H$}：全部为硬——给出两圈系数函数。\\[6pt]
    分组余项：\\[2pt]
    \quad$r=(1-T_H)(1-T_1)(1-T_2)\,\Gamma=\Delta\mathsf H_0\,\Delta K_2\,\Delta K_1$};
  \node[anchor=north west,align=left,font=\footnotesize,text width=6.1cm,draw=red!50!black,line width=0.5pt,
        rounded corners=2pt,fill=red!3,inner sep=4pt] at ($(T.south west)+(0.05,-0.2)$)
    {单个链项 $T_H(1-T_2)\Gamma$ 在 $U_{R_2}$ 上不可积（在 $k_1\parallel\hat p$ 处有无质量共线奇点），因而分组是必要的。};
\end{tikzpicture}
\caption{DIS 两级梯形图（外夸克离壳 $p^2=-m^2$，内线无质量；竖直红色虚线为末态割线，左侧为振幅）的三个区域 $R_2<R_1<R_H$ 与分组余项。
$\Gamma=\mathsf H_0(k_2)\,K_2(k_2,k_1)\,K_1(k_1,p)$，$\mathsf H_0$ 为含两个光子顶点的顶端夸克线，$K_2$、$K_1$ 为上、下两级横档；虚线框依次为 $R_2$、$R_1$、$R_H$ 的硬子图。
分组余项 $r=\Delta\mathsf H_0\Delta K_2\Delta K_1$（$\Delta$ 表示相应因子减去其共线近似）；包含全硬区域 $R_H$ 才能得到两圈系数函数（\cref{ex:disall-ladder}）。}
\label{fig:disall-ladder-regions}
\end{figure}
%%% END fig:disall-ladder-regions

% ---------------------------------------------------------------------
\section{模板假设的验证与主定理}\label{sec:disall-theorem}

\cref{thm:method-template}（全阶论证，[A]）把因子化归结为五个假设：(H1) 领头区域，(H2) 共线近似与 K 胶子的分离，(H3) 软近似与 Glauber 区，(H4) 末态相互作用与软因子，(H5) 减除项与重整化 PDF 的等同。
(H1) 即\cref{prop:disall-regions}（积分形式为\cref{asm:appN-envelope,asm:appN-separation}）；本节验证其余各条并组装主定理。

\begin{proposition}[(H2) 的验证 \tagA]\label{prop:disall-h2}
在 DIS 中，共线近似（\cref{def:method-collinear-approx}）加上 K 胶子的 Ward 恒等式（\cref{lem:method-ward-multi}）与共线 Wilson 线（\cref{prop:method-collinear-wilson}）给出\cref{ch:pdf}的 PDF，
其 Wilson 线是沿 $n_B$ 的直线段：振幅中为未来指向的 $W[\infty n_B,0]$，与共轭振幅中的 $W[w,\infty n_B]$ 合成 $W[w,0]$。色单态靶的色投影为 $\Trc[HJ]\simeq\frac1{N_c}\Trc H\,\Trc J$；
归一化由 $\int\dd^4k=P^+\!\int\dd\xi\int\dd k^-\dd^2\mathbf k_T$ 与 $\int\frac{\dd k^-\dd^2\mathbf k_T}{(2\pi)^4}\Tr[\gamma^+J]=2f$ 给出（\cref{eq:disall-chain}）。（本命题只涉及部分子态。）
\uses{def:method-collinear-approx,lem:method-ward-multi,prop:method-collinear-wilson,prop:method-ward-tree,prop:disall-regions}
% errtype: E-S
\end{proposition}

\begin{proofsketch}
(1) 共线近似作 $H(k)\Phi(k)\to H(\hat k)\frac14\gamma^-\Tr[\gamma^+\Phi(k)]$（对 $\Phi=\sla p$ 给出 $p^+\gamma^-$），其余 Clifford 分量被推动因子压低或对非极化靶为零。
(2) 种子（\cref{prop:method-ward-tree}(b)，\cref{fig:disall-seed}）：$k_\mu H^\mu_a=gt^aH_0$，加上 K 投影即 $W^{(1)}[\infty n_B,0]H_0$；两个 K 胶子精确地给出 $(n_B\cdot k_1)(n_B\cdot k_2)W^{(2)}_{ab}[\infty n_B,0]H_0$，
路径排序来自对易子项；K 胶子附着在末态线上，故 Wilson 线未来指向。(3) 一般阶由\cref{lem:method-ward-multi}与\cref{prop:method-collinear-wilson}：积出 $k^-$、$\mathbf k_T$（\cref{eq:disall-chain}）得到裸 PDF，
$\mathbf w_T=0$ 时 $W[w,\infty n_B]W[\infty n_B,0]=W[w,0]$；色投影来自 $\langle J_{ab}\rangle\propto\delta_{ab}$（\cref{prop:disall-regions}给出的界面结构）。完整论证见\cite[第~11.3~节]{Collins2011}。
\end{proofsketch}

\paragraph{地位与缺口}
最弱的一步是非阿贝尔多胶子 Ward 恒等式（含鬼场）对已被近似子改变过的被积函数成立，即\cref{asm:appN-ward-product}(WI1)，它是\cref{lem:method-ward-multi}（全阶论证，[A]）的精确形式；
阿贝尔情形与树图层次 $n\le2$ 是严格的。该论证中被明确假定的一步是非定域算符的 BRST 分类，\cite[第~11.3.2~节，第~407--410~页]{Collins2011}明言对非定域算符没有证明而予以假定：
即具名缺口\cref{asm:appN-brst-nonlocal}（[C]；只是点名，不是本命题的前提）。全部由纵向胶子诱导硬散射的单个 Feynman 规范图有超领头幂次，对图求和后相消，
但靶中胶子的横动量与规范群的非阿贝尔性留下一个非零的领头项，它正是规范不变胶子 PDF 所需要的\cite{Collins:2008sg}。这一步给出标度陈述（E-S）。

\begin{proposition}[(H3)(H4) 平凡成立 \tagA]\label{prop:disall-h3h4}
在 DIS 中，领头幂次没有软因子，也没有 Glauber 问题：只有一个强子，末态被求和，而非物理区的前向振幅中没有末态喷注。（本命题只涉及部分子态。）
\uses{prop:disall-regions,lem:disall-fsi,lem:method-soft-unitarity}
% errtype: E-S
\end{proposition}

\begin{proofsketch}
(H3)：领头区域没有软子图（\cref{prop:disall-regions}）；Glauber 钳夹需要两个反向运动、不被末态求和消去的共线子图（Drell--Yan 中两个强子的旁观者），而靶残余与流喷注之间的交换在割图求和后不被钳夹（\cref{lem:disall-fsi}）。
(H4)：末态相互作用由\cref{lem:disall-fsi}处理；对共线单举截面，形式上保留的、对割图求和的软因子等于 1（\cref{lem:method-soft-unitarity}，[R]，$R_{\mathrm{def}}$）。
\end{proofsketch}

\paragraph{地位与缺口}
E-S 型论证，强度与\cref{lem:disall-fsi,prop:disall-regions}相同：割图版本依赖于\cref{asm:appN-cutsum}，积分形式依赖于\cref{asm:appN-envelope,asm:appN-separation}。“平凡”只是相对于 Drell--Yan 而言，那里 (H3) 需要附录~O 的 Glauber 分析（\cref{ch:dy}）。

\begin{proposition}[(H5) 的验证 \tagA]\label{prop:disall-h5}
在 DIS 中，减除项等于 $C\otimes f_{\cdot/j}$，其中 $f_{i/j}$ 为重整化的部分子 PDF；因此系数函数由部分子靶算出：$C_{2,i}=\hat F_2^{(j)}\otimes(f^{\mathrm{partonic}})^{-1}_{ji}$（对 $j$ 求和），
在纯维数正规化中单圈为 $C^{(1)}=\hat F^{(1)}+\frac{\as}{2\pi}\frac1\varepsilon P\otimes\hat F^{(0)}$。（本命题只涉及部分子态。）
\uses{prop:pdf-partonic-dr,thm:method-subtraction,prop:disall-factorized-subtraction,asm:appN-multren,asm:appN-collins}
% errtype: E-S
\end{proposition}

\begin{proofsketch}
由\cref{thm:method-subtraction}（[A]，瓶颈\cref{asm:appN-collins}）与\cref{prop:disall-factorized-subtraction}，$T=\sum_i\hat{\mathsf H}^{\mathrm{bare}}_i\otimes f^{\mathrm{bare}}_{i/j}+r$；由\cref{asm:appN-multren}(Z1) 主项为 $Cf^{\mathrm{ren}}$，
$C\equiv\hat{\mathsf H}^{\mathrm{bare}}Z$，由 (Z3)(ii) 它没有极点。对在壳无质量部分子靶，$f^{\mathrm{bare}}_{\cdot/j}=\delta$，$f^{\mathrm{ren}}_{\cdot/j}=Z^{-1}=\delta-\frac{\as}{2\pi}\frac1\varepsilon P+\order{\as^2}$（\cref{prop:pdf-partonic-dr}），
$T^{\mathrm{DR}}_{(j)}=\hat{\mathsf H}^{\mathrm{bare}}$，故 $C=T^{\mathrm{DR}}_{(j)}(f^{\mathrm{ren}}_{\cdot/j})^{-1}$，单圈即所述公式（\cref{sec:disall-nlo}）。
\end{proofsketch}

\paragraph{地位与缺口}
等级由减除条件\cref{asm:appN-collins}与乘法重整化\cref{asm:appN-multren}决定：(Z1) 对 PDF 是定理（壳上规范不变的定域扭度 2 算符），(Z3)(ii) 是维数正规化的质量因子化（[A] 层次，在 NNLO 与 N$^3$LO 系数函数的有限性中被逐阶验证）；
精确路线是\cref{prop:disall-factorized-subtraction}加 (Z3)(ii)。正规子路线\cref{thm:disall-renorm-compat}（内容等级 [C]，经由\cref{asm:appN-ir-regulator}）只是参照而不是依赖：全阶论证的陈述不能以内容等级为 [C] 的前提为依据。

\begin{theorem}[DIS 共线因子化（动量空间，$x$ 空间逐点） \tagA]\label{thm:disall-factorization}
对固定的 $0<x<1$ 且 $W^2=Q^2(1-x)/x+M^2\gg\Lambda^2$，当 $Q\to\infty$ 时
\begin{equation}\label{eq:disall-main}
\frac{F_2(x,Q^2)}{x}=\sum_i\int_x^1\frac{\dd\xi}{\xi}\,C_{2,i}\Big(\frac x\xi,\frac{Q^2}{\mu^2},\as(\mu)\Big)f_i(\xi;\mu)+R(x,Q),
\end{equation}
$F_L$ 同理；其中 $i$ 遍历夸克、反夸克与胶子，$f_i$ 为\cref{def:pdf-quark,def:pdf-gluon}定义的 $\MSbar$ PDF，$C_{2,i}$ 为 $\MSbar$ 减除的部分子结构函数（\cref{prop:disall-h5}）。
余项 $R=\order{\Lambda^2/Q^2}$，其系数依赖于 $x$：当 $x\to1$ 时不一致（启发式的展开参数为 $\Lambda^2/(Q^2(1-x))$），当 $x\to0$ 时亦不一致。
非微扰地位：\cref{asm:intro-bridge}（MA1、MA2、MA4）。
\uses{thm:method-template,prop:disall-regions,lem:disall-fsi,prop:disall-h2,prop:disall-h3h4,prop:disall-h5,prop:disall-two-routes,def:pdf-quark,def:pdf-gluon,asm:intro-bridge,asm:appN-collins,asm:appN-envelope,asm:appN-separation,asm:appN-cutsum,asm:appN-ward-product,asm:appN-uv,asm:appN-multren}
% errtype: E-S
\end{theorem}

\begin{proofsketch}
先对带红外正规化的部分子靶在每一阶证明余项为 $\order{\lambda^aL^K}$，最后过渡到强子。
(1) (H1)：\cref{prop:disall-regions}（标度陈述），积分形式为\cref{asm:appN-envelope,asm:appN-separation}。
(2) (H2)：\cref{prop:disall-h2}，多胶子形式为\cref{asm:appN-ward-product}。
(3) (H3)(H4)：\cref{prop:disall-h3h4}；末态相互作用不被钳夹（\cref{lem:disall-fsi}），割图形式为\cref{asm:appN-cutsum}。
(4) (H5)：\cref{prop:disall-h5}，用到减除定理（瓶颈\cref{asm:appN-collins}）、\cref{asm:appN-uv}与\cref{asm:appN-multren}。
(5) 于是\cref{thm:method-template}的假设全部成立，给出非物理区前向振幅的因子化 $T_a(\omega)=\sum_i\int\frac{\dd\xi}{\xi}\hat C_{a,i}(\xi\omega)f_i(\xi)+\order{\lambda^aL^K}$（$\omega\in[-\rho,\rho]$），
即\cref{prop:disall-two-routes}(a) 与 (H2)、(H5) 的结合。
(6) 回到物理区：主项的 $\omega$ 依赖只在 $\hat C_{a,i}(\xi\omega)$ 中，其右割线在 $\xi\omega\ge1$，故在 $\omega=1/x$ 处的不连续部分只来自 $\xi\ge x$，给出\cref{eq:disall-main}的下限
（对 $\xi<0$，$\hat C$ 在 $\xi\omega\le-1$ 处的交叉割线给出反夸克项，$f_{\bar q}(\xi)=-f_q(-\xi)$）。余项不能这样延拓：复圆盘上的一致界只给出矩（\cref{prop:disall-two-routes}(b)），固定 $x$ 的陈述来自\cref{asm:appN-cutsum}之下对割图求和后的被积函数应用减除定理（\cref{prop:disall-two-routes}(c)），
常数在 $x\to1$ 时增长。这给出部分子形式的\cref{eq:disall-main}。
(7) 强子：\cref{asm:intro-bridge}的 (Br1) 断言没有新的领头区域，(Br2) 断言全部 $\Lambda$ 敏感性都在同一组算符（PDF）的强子矩阵元中，(Br3) 把余项读作 $\order{(\Lambda/Q)^p}$ 的渐近陈述；
对 DIS，$p=2$ 与矩的扭度展开一致，高扭度系数随矩的阶增长，使余项在 $x\to1$ 处不一致，在 $x\to0$ 处则由 $\as\ln(1/x)$ 的增强与饱和效应破坏一致性。
\end{proofsketch}

\paragraph{地位与缺口}
本定理是全阶论证（[A]）。其瓶颈为：\cref{asm:appN-collins}（一般图在积分范数下的减除条件）、\cref{asm:appN-envelope}与\cref{asm:appN-separation}（幂次计数的积分形式）、
\cref{asm:appN-cutsum}（割图求和的公共表示）、\cref{asm:appN-ward-product}（多胶子 Ward 恒等式的乘积结构）、\cref{asm:appN-uv}（被积函数层次的紫外重整化）、
\cref{asm:appN-multren}（共线因子的乘法重整化），以及强子层次的\cref{asm:intro-bridge}。文献给出这些步骤的标度陈述（E-S）；在附录~N 的编号假设之下余项在积分范数意义下（E-N）受控，条件形式为\cref{thm:appN-structure}。
该论证中被明确假定的一步是非定域算符的 BRST 分类（\cite[第~11.3.2~节，第~407--410~页]{Collins2011}），即具名缺口\cref{asm:appN-brst-nonlocal}（[C]），它经由\cref{lem:method-ward-multi}与\cref{thm:method-template}进入，
只是点名，不是本定理的前提；朴素的 K 胶子恒等式还有非零的领头余项\cite{Collins:2008sg}。非规范理论中的完整证明见\cite[第~8.9~节]{Collins2011}；对 QCD，\cite[第~10.1~节]{CollinsSoperSterman1989}写道
「“there is as yet in the literature no complete and systematic subtraction procedure in QCD \ldots{} even in the case of deeply inelastic scattering”」，即即使对 DIS，文献中也还没有完整而系统的 QCD 减除程序。
早期的全阶论证\cite{EllisEtAl1979,Ellis:1978sf,AmatiPetronzioVeneziano1978a,AmatiPetronzioVeneziano1978b,Mueller1978,CurciFurmanskiPetronzio1980}多在物理规范（轴规范、光锥规范）中进行；
物理规范的非物理奇点阻碍把围道变形出 Glauber 区，因而一般而可靠的证明应在 Feynman 规范中进行\cite[第~12.6~节]{Collins2011}。提升本定理须解除上述假设；$x\to1$ 处的一致性还需要阈值区的重新因子化（\cref{ch:thr}）。

% ---------------------------------------------------------------------
\section{与 OPE 结果的一致性}\label{sec:disall-consistency}

OPE 的 Wilson 系数 $C^{(n)}_{a,i}$（\cref{ch:dis}）与动量空间证明中由部分子靶构造的系数函数（\cref{prop:disall-h5}）必须一致；一致的是作为分布的扭度 2 部分，而不是逐点的余项。

\begin{corollary}[与 OPE 系数的一致性 \tagR]\label{cor:disall-consistency}
（$R_{\mathrm{OPE}}$：以 OPE 为前提而严格。QCD 中含夸克的微扰 OPE 是全阶论证（\cref{thm:dis-ope}），严格证明只对标量理论与纯 Yang--Mills 的规范不变算符给出；对强子态还须非微扰 OPE（\cref{asm:intro-np-ope}）。）
以\cref{prop:pdf-partonic-dr}的构造定义动量空间的系数函数 $C^{\mathrm{mom}}_{a,i}$：在纯维数正规化中对在壳无质量部分子靶 $j$，$\hat F^{(j)}_a/z=\sum_iC^{\mathrm{mom}}_{a,i}\otimes f_{i/j}$，$f_{i/j}$ 为 $\MSbar$ 部分子 PDF。
则对 OPE 有定义的一切 $n$，$\int_0^1\dd z\,z^{n-1}C^{\mathrm{mom}}_{a,i}(z)=C^{(n)}_{a,i}$；因此 $C^{\mathrm{mom}}_{a,i}$ 就是\cref{thm:pdf-dis-xspace}中以其矩唯一确定的分布 $C_{a,i}$，
\cref{thm:disall-factorization}与\cref{thm:pdf-dis-xspace}的扭度 2 部分作为分布相同；逐点的余项只受全阶论证\cref{thm:disall-factorization}控制。（本推论只涉及部分子态的系数函数。）
\uses{thm:pdf-dis-xspace,prop:dis-wilson-partons,cor:appK-pdf,prop:pdf-partonic-dr,lem:prelim-mellin,thm:pdf-moments-local}
\end{corollary}

\begin{proof}
(1) 由 Mellin 卷积定理（\cref{lem:prelim-mellin}），$\int_0^1\dd z\,z^{n-1}(C\otimes f)(z)=C_nf_n$，$X_n\equiv\int_0^1\dd z\,z^{n-1}X(z)$。对定义式取矩：
$\int_0^1\dd z\,z^{n-2}\hat F^{(j)}_2=\sum_iC^{\mathrm{mom}}_{2,i,n}\,f_{i/j,n}$。
(2) OPE 一侧，\cref{prop:dis-wilson-partons}（以 OPE 为前提而严格）给出 $\int_0^1\dd z\,z^{n-2}\hat F^{(j)}_2=\sum_iC^{(n)}_{2,i}\langle j|O^{(n)}_i|j\rangle$，
其中在壳无质量部分子的算符矩阵元的全部圈修正都无标度，裸矩阵元等于 $\delta_{ij}$，重整化后等于 $\delta_{ij}$ 加纯 $\MSbar$ 极点。由\cref{thm:pdf-moments-local}，
重整化的扭度 2 定域算符矩阵元等于重整化 PDF 的矩；对部分子靶，二者都等于“裸值 $\delta_{ij}$ 乘以同一个 $Z^{-1}$”（\cref{prop:pdf-partonic-dr}），故 $\langle j|O^{(n)}_i|j\rangle=f_{i/j,n}$。
(3) 两式相减：$\sum_i\big(C^{\mathrm{mom}}_{2,i,n}-C^{(n)}_{2,i}\big)f_{i/j,n}=0$ 对一切部分子靶 $j$ 成立；矩阵 $(f_{i/j,n})=\mathbb 1+\order{\as}$ 逐阶可逆，故逐阶地 $C^{\mathrm{mom}}_{2,i,n}=C^{(n)}_{2,i}$。$F_L$（以及 $F_1$）同理。
(4) $C^{\mathrm{mom}}_{2,i}$ 是 $[0,1]$ 上的有限阶分布（$\delta(1-z)$、加号分布与可积函数）；由\cref{cor:appK-pdf}的论证，偶数阶矩（$n\ge N_1$）唯一地确定这样的分布，
所以 $C^{\mathrm{mom}}_{2,i}$ 等于\cref{thm:pdf-dis-xspace}中矩为 $C^{(n)}_{2,i}$ 的那个分布。于是两个定理中 $x$ 空间的主项逐项相同。
(5) 反过来，矩的相同不能推出余项逐点相同：对每个固定 $n$ 矩为 $\order{1/Q^2}$ 的函数，逐点可以是 $\order{1}$，例如 $\sin(Q^2x)$。所以\cref{thm:pdf-dis-xspace}只控制分布意义下的余项，
固定 $x$ 处的逐点余项是全阶论证\cref{thm:disall-factorization}的内容。
\end{proof}

\noindent
在 $\order{\as}$ 这一一致性已被逐项检验：在壳部分子前向 Compton 振幅的硬区展开，其 Taylor 系数作为 $\varepsilon$ 的有理函数与割图结果之矩精确相等（夸克靶 $N\le8$，胶子靶偶数 $N\le8$），
重整化后给出\cref{ch:dis}的闭式（\cref{ex:dis-oneloop-ope}）；下一节 $x$ 空间系数函数的矩也与之一致。   % VF-DIS-13
逐点余项、$x\to1$ 与共振区（$W^2\sim\Lambda^2$，只能在对 $x$ 平滑后与扭度展开比较，Bloom--Gilman 对偶，猜想层次）的情形见\cref{sec:disall-summary}的常见误解。

% ---------------------------------------------------------------------
\section{$\order{\as}$ 系数函数的推导}\label{sec:disall-nlo}

本节在 $D=4-2\varepsilon$ 维的常规维数正规化中计算夸克靶与胶子靶的部分子结构函数至 $\order{\as}$（图见\cref{fig:dis-oneloop}），并由 $\MSbar$ 质量因子化得到系数函数：
软—共线双极点与末态共线极点在实、虚修正之和中相消，剩下的初态共线极点恰好由 $\MSbar$ 部分子 PDF 吸收——这正是\cref{thm:disall-factorization}的机制在单圈的实现。

\begin{proposition}[$O(\alpha_s)$ 的 $\overline{\rm MS}$ 系数函数 \tagR]\label{prop:disall-nlo}
（$R_{\mathrm{PT}}$：在重整化微扰论中对红外正规化的部分子态严格。）以下陈述在单圈（$\order{\as}$）。（本命题只涉及部分子态。）
在常规维数正规化、入射胶子按 $D-2$ 个极化平均的 $\MSbar$ 方案中，按
\begin{equation}\label{eq:disall-norm}
\frac{F_2}{x}=\sum_qe_q^2\Big\{q+\bar q+\frac{\as}{2\pi}\big[C_2^q\otimes(q+\bar q)+2C_2^g\otimes g\big]\Big\},\qquad
\frac{F_L}{x}=\sum_qe_q^2\frac{\as}{2\pi}\big[C_L^q\otimes(q+\bar q)+2C_L^g\otimes g\big]
\end{equation}
归一化系数函数（$\sum_q$ 只对味求和；胶子项的因子 $2$ 计入 $\gamma^*g\to q\bar q$ 中同时出现的夸克与反夸克；$(C\otimes f)(x)=\int_x^1\frac{\dd\xi}\xi C(x/\xi)f(\xi)$），在 $\mu=Q$ 处
\begin{equation}\label{eq:disall-C}
\begin{aligned}
C_2^q(z)&=C_F\Big[2\mathcal D_1-\frac32\mathcal D_0-(1+z)\ln(1-z)-\frac{1+z^2}{1-z}\ln z+3+2z-\Big(\frac92+\frac{\pi^2}3\Big)\delta(1-z)\Big],\\
C_2^g(z)&=T_F\Big[\big(z^2+(1-z)^2\big)\ln\frac{1-z}{z}-8z^2+8z-1\Big],\\
C_L^q(z)&=2C_Fz,\qquad C_L^g(z)=4T_Fz(1-z),
\end{aligned}   % VF-DIS-7
\end{equation}
$\mathcal D_k\equiv[\ln^k(1-z)/(1-z)]_+$。对一般的 $\mu=\mu_F=\mu_R$，$C_2^i(z,Q^2/\mu^2)=C_2^i(z)+\ln(Q^2/\mu^2)P_{qi}(z)$（$i=q,g$），$C_L$ 在此阶与 $\mu$ 无关。
等价地 $C_2^q=C_F\big[\frac{1+z^2}{1-z}\big(\ln\frac{1-z}z-\frac34\big)+\frac{9+5z}4\big]_+$（加号作用于整个方括号）。若入射胶子按 2 个极化平均，则 $C_2^g\to C_2^g+P_{qg}$，$C_L^g$ 不变。
\uses{def:disall-breit,prop:sudakov-oneloop,lem:prelim-plus-expansion,prop:pdf-partonic-dr,lem:appC-phase-space,prop:dis-c-oneloop}
\end{proposition}

\begin{proof}
部分子结构函数 $\hat F^j_a/(ze_q^2)$（$j=q,g$）用\cref{def:disall-breit}的投影提取；$g_0^2=4\pi\as\mu^{2\varepsilon}S_\varepsilon^{-1}$，$S_\varepsilon=(4\pi)^\varepsilon\ee^{-\varepsilon\gamma_E}$；初态平均为夸克自旋 $\frac12$、颜色 $\frac1{N_c}$，胶子极化 $\frac1{D-2}$、颜色 $\frac1{N_c^2-1}$。
(1) Born：$\hat W^{\mu\nu}_{(0)}=\frac1{4\pi}2\pi\delta((p+q)^2)\frac12e_q^2\Tr[\sla p\gamma^\mu(\sla p+\sla q)\gamma^\nu]$；由 $\gamma^\mu\sla a\gamma_\mu=-(D-2)\sla a$、$\sla p\sla p=0$ 与
$\delta((p+q)^2)=\delta(1-z)/Q^2$，对任意 $D$ 有 $\hat F^{q(0)}_2/z=e_q^2\delta(1-z)$，$\hat F^{q(0)}_L=0$。   % VF-DIS-1
(2) 虚修正：由一圈形状因子（\cref{prop:sudakov-oneloop}），$\hat F^{q,\mathrm V}_2/(ze_q^2)=2F^{(1)}\delta(1-z)=\frac{\as}{2\pi}C_F(\frac{\mu^2}{Q^2})^\varepsilon[-\frac2{\varepsilon^2}-\frac3\varepsilon-8+\zeta_2+\order{\varepsilon}]\delta(1-z)$，极点全部是红外的（$Z_1=Z_2$，外线自能无标度），对 $F_L$ 无贡献。   % VF-DIS-3
(3) 实辐射 $\gamma^*q\to qg$：以 $D$ 维两体相空间（\cref{lem:appC-phase-space}）积分投影后的平方矩阵元，$y$ 积分均为 Beta 函数，对 $\varepsilon$ 精确地得到（中间步骤见\cref{ex:disall-nlo}）
\begin{equation}\label{eq:disall-Rq}
\begin{gathered}
\frac{\hat F^{q,\mathrm R}_a}{ze_q^2}=\frac{\as}{2\pi}C_F\Big(\frac{\mu^2}{Q^2}\Big)^\varepsilon R^q_a,\\
R^q_2=\frac{\ee^{\varepsilon\gamma_E}\Gamma(-\varepsilon)}{2\Gamma(2-2\varepsilon)}\frac{z^\varepsilon\,[2(1+z^2)-\varepsilon(7-2z)+\varepsilon^2]}{(1-z)^{1+\varepsilon}},\qquad
R^q_L=\frac{\ee^{\varepsilon\gamma_E}\Gamma(2-\varepsilon)}{\Gamma(2-2\varepsilon)}\,2z\Big(\frac z{1-z}\Big)^\varepsilon .   % VF-DIS-4
\end{gathered}
\end{equation}
(4) 用分布展开（\cref{lem:prelim-plus-expansion}，$a=1$）$(1-z)^{-1-\varepsilon}=-\frac1\varepsilon\delta(1-z)+\sum_{k\ge0}\frac{(-\varepsilon)^k}{k!}\mathcal D_k$ 与乘积约定，记 $\delta\equiv\delta(1-z)$，
\begin{equation}\label{eq:disall-Rq-exp}
\begin{aligned}
R^q_2={}&\frac{2\delta}{\varepsilon^2}+\frac1\varepsilon\Big[-2\mathcal D_0+1+z+\frac32\delta\Big]+2\mathcal D_1-\frac32\mathcal D_0\\
&-(1+z)\ln(1-z)-\frac{1+z^2}{1-z}\ln z+3+2z+\Big(\frac72-3\zeta_2\Big)\delta+\order{\varepsilon},   % VF-DIS-4
\end{aligned}
\end{equation}
而 $R^q_L=2z+\order{\varepsilon}$ 有限。(5) 实、虚相加：$1/\varepsilon^2$ 的系数 $2-2=0$；$1/\varepsilon$ 的系数 $-2\mathcal D_0+1+z+(\frac32-3)\delta=-P_{qq}/C_F$（$P_{qq}/C_F=2\mathcal D_0-(1+z)+\frac32\delta$）；
$\delta$ 的有限系数 $\frac72-3\zeta_2-8+\zeta_2=-(\frac92+\frac{\pi^2}3)$。于是 $\hat F^q_2/(ze_q^2)=\delta+\frac{\as}{2\pi}(\frac{\mu^2}{Q^2})^\varepsilon[-\frac1\varepsilon P_{qq}+C_2^q+\order{\varepsilon}]$。   % VF-DIS-6
(6) 质量因子化：由\cref{prop:pdf-partonic-dr}，$f_{q/q}=\delta-\frac{\as}{2\pi}\frac1\varepsilon P_{qq}+\order{\as^2}$；因子化公式用于夸克靶给出
$\frac{\as}{2\pi}C_2^q(z,Q^2/\mu^2)=\frac{\as}{2\pi}[(\frac{\mu^2}{Q^2})^\varepsilon(-\frac{P_{qq}}\varepsilon+C_2^q)+\frac{P_{qq}}\varepsilon]+\order{\varepsilon}=\frac{\as}{2\pi}[C_2^q+\ln\frac{Q^2}{\mu^2}P_{qq}]$；
胶子靶的极点由 $f_{q/g}+f_{\bar q/g}=-\frac{\as}{2\pi}\frac2\varepsilon P_{qg}$ 吸收（\cref{ex:disall-gluon}）。(7) 两种写法的 $C_2^q$ 在 $z<1$ 处逐点相同，且在 $[0,1]$ 上积分都为零（Adler 求和规则），故为同一分布；
按 2 个极化平均相当于把 $R^g_2$ 乘以 $1-\varepsilon$，极点项 $-\frac2\varepsilon P_{qg}$ 产生有限移动 $+2P_{qg}$，即 $C_2^g\to C_2^g+P_{qg}$。   % VF-CONV-6
这些结果的 Mellin 矩就是\cref{ch:dis}的单圈 Wilson 系数（\cref{prop:dis-c-oneloop}）。
\end{proof}

\noindent
换到 $a_s=\as/4\pi$ 的归一化（$c^{(1)}_{2,\mathrm{ns}}=2C_2^q$，$c^{(1)}_{2,g}=4n_fC_2^g$，$c^{(1)}_{L,\mathrm{ns}}=2C_L^q$，$c^{(1)}_{L,g}=4n_fC_L^g$，$T_F=\frac12$），上述结果与\cite[式~(4.1)--(4.4)、(B.3)--(B.4)]{Vermaseren:2005qc}
及\cite[式~(3)]{Moch:2004xu}逐项一致，那里归功于\cite{BardeenEtAl1978}（其式~(3.9) 去掉 $\ln4\pi-\gamma_E$ 项后恰为 $2C^q_{2,N}$）；参见\cite{FurmanskiPetronzio1982,EllisStirlingWebber1996}。

\begin{example}[$O(\alpha_s)$ 系数函数的完整推导]\label{ex:disall-nlo}
我们写出\cref{prop:disall-nlo}证明第 (3)(4) 步的中间计算。实辐射振幅（\cref{fig:dis-oneloop}）为
$M^\mu=e_qgt^a\bar u(p')[\gamma^\rho\frac{\sla p+\sla q}{\hat s}\gamma^\mu+\gamma^\mu\frac{\sla p-\sla k}{\hat t}\gamma^\rho]u(p)\varepsilon^*_\rho(k)$，$\hat s=Q^2(1-z)/z$，$\hat t=(p-k)^2$，$\hat u=(p-p')^2$，$\hat s+\hat t+\hat u=-Q^2$。
只有一个外胶子时 $k^\rho M_\rho=0$，极化求和可取 $-g_{\rho\sigma}$（我们也用物理极化求和验证了结果不变）。$D$ 维 Dirac 代数给出
$(-g_{\mu\nu})\sum M^\mu M^{\nu*}|_{\mathrm{Dirac}}=8(1-\varepsilon)[(1-\varepsilon)(-\frac{\hat t}{\hat s}-\frac{\hat s}{\hat t})+\frac{2\hat uQ^2}{\hat s\hat t}+2\varepsilon]$，$(p_\mu p_\nu)\sum M^\mu M^{\nu*}|_{\mathrm{Dirac}}=-4(1-\varepsilon)\hat u$，
再乘以 $e_q^2g_0^2$ 与色—自旋平均因子 $\frac12\cdot\frac1{N_c}\Tr(t^at^a)=\frac{C_F}2$。在 $\gamma^*$—部分子质心系中以 $y=(1-\cos\theta)/2$ 参数化（$\theta$ 为出射胶子与 $\vec p$ 的夹角），
$\hat t=-Q^2y/z$，$\hat u=-Q^2(1-y)/z$，$\int\dd\Phi_2=\frac1{8\pi}(\frac{4\pi}{\hat s})^\varepsilon\frac1{\Gamma(1-\varepsilon)}\int_0^1\dd y\,[y(1-y)]^{-\varepsilon}$，总前因子为
\begin{equation*}
\frac1{4\pi}\cdot\frac{C_F}2g_0^2\cdot\frac1{8\pi}\Big(\frac{4\pi}{\hat s}\Big)^\varepsilon\frac1{\Gamma(1-\varepsilon)}=\frac{\as}{2\pi}C_F\cdot\frac18\Big(\frac{\mu^2}{Q^2}\Big)^\varepsilon\frac{z^\varepsilon(1-z)^{-\varepsilon}\ee^{\varepsilon\gamma_E}}{\Gamma(1-\varepsilon)}.
\end{equation*}
$F_2$ 投影后的被积函数是 $8[(1-\varepsilon)y^2-2y(z-\varepsilon(1-z))+1+z^2-\varepsilon(1-z)^2]/[y(1-z)]+16(3-2\varepsilon)z(1-y)$。
记 $I_a=\int_0^1\dd y\,y^{a-\varepsilon}(1-y)^{-\varepsilon}$，则 $I_0=\Gamma^2(1-\varepsilon)/\Gamma(2-2\varepsilon)$，$I_1=I_0/2$，$I_{-1}=-(1-2\varepsilon)I_0/\varepsilon$；$y\to0$ 端点（胶子与入射夸克共线）的 $1/\varepsilon$ 来自 $I_{-1}$。
逐项积分并乘以 $(1-z)/(8I_0)$ 得 $-(1+z^2)/\varepsilon+(7-2z)/2-\varepsilon/2=-\frac1{2\varepsilon}[2(1+z^2)-\varepsilon(7-2z)+\varepsilon^2]$，连同前因子与 $\Gamma(-\varepsilon)=-\Gamma(1-\varepsilon)/\varepsilon$ 即得\cref{eq:disall-Rq}。
$F_L$ 投影后的被积函数 $32(1-\varepsilon)z(1-y)$ 在 $y\to0$ 处不奇异：$F_L$ 的投影压低了共线区域，$R^q_L$ 有限。

展开时写 $R^q_2=A(\varepsilon)(1-z)^{-1-\varepsilon}G(z,\varepsilon)$，$A=\frac{\ee^{\varepsilon\gamma_E}\Gamma(-\varepsilon)}{2\Gamma(2-2\varepsilon)}$，$G=z^\varepsilon[2(1+z^2)-\varepsilon(7-2z)+\varepsilon^2]$，$G(1,\varepsilon)=4-5\varepsilon+\varepsilon^2$。
于是 $\delta$ 的系数为 $-AG(1,\varepsilon)/\varepsilon=\frac2{\varepsilon^2}+\frac3{2\varepsilon}+\frac72-3\zeta_2+\order{\varepsilon}$，$\mathcal D_0$ 的系数为 $AG(1,\varepsilon)=-\frac2\varepsilon-\frac32+\order{\varepsilon}$，$\mathcal D_1$ 的系数为 $-\varepsilon AG(1,\varepsilon)=2+\order{\varepsilon}$，
正规部分 $A[G(z,\varepsilon)-G(1,\varepsilon)](1-z)^{-1-\varepsilon}=\frac{1+z}\varepsilon-(1+z)\ln(1-z)-\frac{1+z^2}{1-z}\ln z+3+2z+\order{\varepsilon}$，合起来即\cref{eq:disall-Rq-exp}。
最后，对 $N=1,\dots,8$ 以 40 位精度数值计算\cref{eq:disall-C}的 Mellin 矩，与\cref{ch:dis}的闭式一致到 $10^{-39}$；特别地 $C^q_{2,1}=0$（Adler 求和规则），   % VF-DIS-8
$C^q_{1,1}=C^q_{2,1}-C^q_{L,1}=-C_F$（Bjorken $\nu$ 求和规则；$C_1^i\equiv C_2^i-C_L^i$ 来自 $2F_1=F_2/x-F_L/x$）。
\uses{prop:disall-nlo}
\end{example}

\begin{example}[光子—胶子融合]\label{ex:disall-gluon}
$\gamma^*(q)+g(p)\to q(k_1)+\bar q(k_2)$：$M^\mu=e_qgt^a\bar u(k_1)[\gamma^\mu\frac{\sla k_1-\sla q}{\hat u}\gamma^\rho+\gamma^\rho\frac{\sla k_1-\sla p}{\hat t}\gamma^\mu]v(k_2)\varepsilon_\rho(p)$，
$\hat t=(p-k_1)^2=-Q^2y/z$，$\hat u=(p-k_2)^2=-Q^2(1-y)/z$（$\theta$ 为出射夸克与 $\vec p$ 的夹角；矩阵元在 $\hat t\leftrightarrow\hat u$ 下对称，\cref{fig:dis-oneloop}中按内线动量对两个传播子的命名不影响结果）。
投影后的迹为 $(-g_{\mu\nu})\sum M^\mu M^{\nu*}|_{\mathrm{Dirac}}=8(1-\varepsilon)[(1-\varepsilon)(\frac{\hat u}{\hat t}+\frac{\hat t}{\hat u})-\frac{2\hat sQ^2}{\hat t\hat u}-2\varepsilon]$，$(p_\mu p_\nu)\sum M^\mu M^{\nu*}|_{\mathrm{Dirac}}=8\hat s$，
平均因子为 $\frac1{2(1-\varepsilon)}\frac1{N_c^2-1}\Tr(t^at^a)=\frac{T_F}{2(1-\varepsilon)}$。记 $\hat F^g_a/(ze_q^2)=\frac{\as}{2\pi}T_F(\frac{\mu^2}{Q^2})^\varepsilon R^g_a$（一味夸克，$q$ 与 $\bar q$ 都计入），
$F_2$ 投影的 $y$ 积分为 $-\frac{16}\varepsilon\frac{\Gamma^2(1-\varepsilon)}{\Gamma(2-2\varepsilon)}\frac{P-3\varepsilon+4\varepsilon^2-2\varepsilon^3}{1-\varepsilon}$，$P\equiv z^2+(1-z)^2$，于是
\begin{equation*}
R^g_2=\frac{2\ee^{\varepsilon\gamma_E}\Gamma(-\varepsilon)}{\Gamma(2-2\varepsilon)}\Big(\frac z{1-z}\Big)^\varepsilon\frac{P-3\varepsilon+4\varepsilon^2-2\varepsilon^3}{(1-\varepsilon)^2}
=-\frac2\varepsilon P+2\Big\{P\ln\frac{1-z}z-8z^2+8z-1\Big\}+\order{\varepsilon},   % VF-DIS-5
\end{equation*}
而 $R^g_L=8z(1-z)+\order{\varepsilon}$。极点出现两次，$y\to0$ 与 $y\to1$（夸克或反夸克与胶子共线）各贡献 $-P/\varepsilon$。
因此 $\hat F^g_2/(ze_q^2)=\frac{\as}{2\pi}(\frac{\mu^2}{Q^2})^\varepsilon[-\frac2\varepsilon P_{qg}+2C_2^g+\order{\varepsilon}]$，减去 $f_{q/g}+f_{\bar q/g}=-\frac{\as}{2\pi}\frac2\varepsilon P_{qg}$ 与 Born 系数之卷积即得 $2[C_2^g+\ln(Q^2/\mu^2)P_{qg}]$。
按 2 个极化平均时 $R^g_2$ 乘以 $1-\varepsilon$，$C_2^g$ 移动 $+P_{qg}$：胶子平均的约定是 $\MSbar$ 方案的一部分。
\uses{prop:disall-nlo}
\end{example}

\noindent
Callan--Gross 关系 $F_L=0$ 只是部分子模型（领头阶）的结果，而不是 QCD 的定理：由 $C_L^q=2C_Fz$、$C_L^g=4T_Fz(1-z)$ 与领头阶的 $F_2/x=\sum e_q^2(q+\bar q)$ 立即得到
$F_L=\frac{\as}{4\pi}x^2\int_x^1\frac{\dd\xi}{\xi^3}\big[\frac{16}3F_2(\xi)+8\sum_qe_q^2(1-\frac x\xi)\xi g(\xi)\big]$（$C_F=\frac43$，$T_F=\frac12$），即 Altarelli--Martinelli 关系\cite{AltarelliMartinelli1978}；
$C_L^q$ 的矩最早由\cite{Zee:1974du}得到。$C_L$ 在 $\order{\as}$ 与 $\mu$ 及方案都无关，是胶子分布的直接检验之一。

% ---------------------------------------------------------------------
\section{标度无关性与 DIS 方案}\label{sec:disall-scale}

$\mu$ 是把共线发散分给 PDF 时引入的人为标度。DGLAP 方程（\cref{thm:pdf-dglap}）是算符重整化的结论，逻辑上先于因子化定理；标度无关性是它与系数函数的 $\mu$ 依赖相互抵消的推论。

\begin{proposition}[标度无关性 \tagR]\label{prop:disall-scale}
（$R_{\mathrm{PT}}$：在重整化微扰论中对红外正规化的部分子态严格。）以下陈述针对截断到 $\order{\as}$ 的\cref{eq:disall-norm}。（本命题只涉及部分子态的系数函数。）
设 PDF 满足 DGLAP 方程，$a\equiv\as(\mu_R)/2\pi$ 满足 $\dd a/\dd\ln\mu_R^2=-\frac{\beta_0}2a^2+\order{a^3}$，系数函数取\cref{prop:disall-nlo}的 $C_2^i(z,Q^2/\mu_F^2)$。则
\begin{equation}\label{eq:disall-scale}
\frac{\partial}{\partial\ln\mu_F^2}\frac{F_2}x=\order{a^2},\qquad
\frac{\partial}{\partial\ln\mu_R^2}\frac{F_2}x=-\frac{\beta_0}2a^2\sum_qe_q^2\big[C_2^q\otimes(q+\bar q)+2C_2^g\otimes g\big]+\order{a^3},   % VF-DIS-10 (written for F_2/x)
\end{equation}
式中 $C_2^i$ 即 $C_2^i(z,Q^2/\mu_F^2)$；且 $F_L$ 在 $\order{a}$ 与 $\mu_F$、$\mu_R$ 都无关。
\uses{prop:disall-nlo,thm:pdf-dglap,prop:prelim-beta}
\end{proposition}

\begin{proof}
在 Mellin 空间中卷积成为乘积。对一味夸克，记 $\Sigma_N$、$g_N$ 为 $q+\bar q$ 与 $g$ 的矩，$F_{2,N}$ 为 $F_2/x$ 的矩，$L_F=\ln(Q^2/\mu_F^2)$，则
$F_{2,N}=e_q^2\{\Sigma_N(\mu_F)+a(\mu_R)[(C^q_{2,N}+L_FP_{qq,N})\Sigma_N+2(C^g_{2,N}+L_FP_{qg,N})g_N]\}$。由\cref{thm:pdf-dglap}，
$\partial\Sigma_N/\partial\ln\mu_F^2=a(\mu_F)[P_{qq,N}\Sigma_N+2P_{qg,N}g_N]+\order{a^2}$（$q$ 与 $\bar q$ 各从胶子得到一份 $P_{qg}\otimes g$），而 $\partial L_F/\partial\ln\mu_F^2=-1$，故
$\partial F_{2,N}/\partial\ln\mu_F^2=e_q^2[a(\mu_F)-a(\mu_R)](P_{qq,N}\Sigma_N+2P_{qg,N}g_N)+\order{a^2}=\order{a^2}$。对 $\mu_R$，只有 $a(\mu_R)$ 依赖于它，
由\cref{prop:prelim-beta}，$\dd a/\dd\ln\mu^2=\frac1{2\pi}\dd\as/\dd\ln\mu^2=-\frac{\beta_0}{8\pi^2}\as^2=-\frac{\beta_0}2a^2$，即得第二式。
$F_L$ 在 $\order{a}$ 只含与 $\mu$ 无关的 $C_L$ 乘以 PDF，而 PDF 的 $\mu$ 依赖是 $\order{a}$，故 $F_L$ 的变化为 $\order{a^2}$。
（以上用符号计算在 Mellin 空间中核对过。）
\end{proof}

\noindent
定理只保证剩余的标度依赖是高一阶的；把 $\mu$ 在 $Q/2$ 与 $2Q$ 之间变化的差别当作“理论误差”是没有定理支持的约定。
另一种常用方案把 $F_2$ 的全部高阶修正吸收进夸克分布；方案变换的一般形式是 $f'=Z\otimes f$、$C'=C\otimes Z^{-1}$（\cref{rem:pdf-schemes}）。

\begin{proposition}[DIS 方案 \tagR]\label{prop:disall-dis-scheme}
（$R_{\mathrm{PT}}$：在重整化微扰论中对红外正规化的部分子态严格。）以下陈述在 $\order{\as}$。（本命题只涉及部分子态的系数函数与方案变换。）
DIS 方案要求 $F_2/x=\sum_qe_q^2(q+\bar q)^{\mathrm{DIS}}$ 在各阶成立。在 $\order{\as}$（$a=\as/2\pi$，全部量取在同一 $\mu$，$C_2^i\equiv C_2^i(z)$ 为 $\mu=Q$ 处的 $\MSbar$ 系数），
\begin{equation}\label{eq:disall-dis-scheme}
\begin{gathered}
q_i^{\mathrm{DIS}}=q_i+a\big[C_2^q\otimes q_i+C_2^g\otimes g\big],\qquad \bar q_i^{\mathrm{DIS}}=\bar q_i+a\big[C_2^q\otimes\bar q_i+C_2^g\otimes g\big],\\
g^{\mathrm{DIS}}=g-a\big[C_2^q\otimes\Sigma+2n_fC_2^g\otimes g\big],   % VF-DIS-11
\end{gathered}
\end{equation}
$\Sigma=\sum_i(q_i+\bar q_i)$；胶子的变换不由 $F_2$ 确定，上式是对全部矩保持动量求和规则的选择。于是 $C_2^{q,\mathrm{DIS}}=\ln(Q^2/\mu^2)P_{qq}$，$C_2^{g,\mathrm{DIS}}=\ln(Q^2/\mu^2)P_{qg}$；
领头阶分裂函数不变，次领头阶在 $(\Sigma,g)$ 基中按
\begin{equation}\label{eq:disall-P1}
\begin{gathered}
P'^{(1)}=P^{(1)}+Z\otimes P^{(0)}-P^{(0)}\otimes Z-\frac{\beta_0}2Z,\qquad Z=\begin{pmatrix}C_2^q&2n_fC_2^g\\-C_2^q&-2n_fC_2^g\end{pmatrix},\\
P^{(1),\mathrm{DIS}}_{\mathrm{NS}}=P^{(1),\MSbar}_{\mathrm{NS}}-\frac{\beta_0}2C_2^q   % VF-DIS-11
\end{gathered}
\end{equation}
变换（DGLAP 核写作 $aP^{(0)}+a^2P^{(1)}$）；$C_L$ 在 $\order{\as}$ 与方案无关。
\uses{prop:disall-nlo,rem:pdf-schemes,thm:pdf-dglap,prop:prelim-beta}
\end{proposition}

\begin{proof}
(1) 由\cref{eq:disall-dis-scheme}，$q_i=q_i^{\mathrm{DIS}}-a[C_2^q\otimes q_i+C_2^g\otimes g]+\order{a^2}$。代入\cref{eq:disall-norm}（取一般 $\mu$，$C_2^i(z,Q^2/\mu^2)=C_2^i(z)+\ln(Q^2/\mu^2)P_{qi}$），
$\order{a}$ 中的 $C_2^q\otimes(q+\bar q)+2C_2^g\otimes g$ 恰好抵消，剩下
$F_2/x=\sum_qe_q^2\{(q+\bar q)^{\mathrm{DIS}}+a\ln(Q^2/\mu^2)[P_{qq}\otimes(q+\bar q)+2P_{qg}\otimes g]\}+\order{a^2}$，即 $C_2^{i,\mathrm{DIS}}=\ln(Q^2/\mu^2)P_{qi}$，在 $\mu=Q$ 处为零。
(2) $\Sigma^{\mathrm{DIS}}+g^{\mathrm{DIS}}=\Sigma+g+a[C_2^q\otimes\Sigma+2n_fC_2^g\otimes g]-a[C_2^q\otimes\Sigma+2n_fC_2^g\otimes g]=\Sigma+g$，对每个矩都成立，特别地 $N=2$ 的动量求和规则不变。
(3) $F_L$ 是 $\order{a}$ 的量，PDF 的 $\order{a}$ 变化只改变 $\order{a^2}$ 项，所以 $C_L$ 不变；领头阶分裂函数也不变。
(4) 设 $f'=(1+aZ)\otimes f$，$\dd f/\dd\ln\mu^2=(aP^{(0)}+a^2P^{(1)})\otimes f$，$\dd a/\dd\ln\mu^2=-\frac{\beta_0}2a^2$（\cref{thm:pdf-dglap,prop:prelim-beta}）。则
$\dd f'/\dd\ln\mu^2=(1+aZ)\otimes(aP^{(0)}+a^2P^{(1)})\otimes f-\frac{\beta_0}2a^2Z\otimes f$；用 $f=(1-aZ)\otimes f'+\order{a^2}$ 消去 $f$，得
$\dd f'/\dd\ln\mu^2=aP^{(0)}\otimes f'+a^2[P^{(1)}+Z\otimes P^{(0)}-P^{(0)}\otimes Z-\frac{\beta_0}2Z]\otimes f'+\order{a^3}$。在 $(\Sigma,g)$ 基中 $Z$ 的两行分别来自 $\Sigma^{\mathrm{DIS}}$ 与 $g^{\mathrm{DIS}}$ 的定义。
非单态部分 $Z_{\mathrm{NS}}=C_2^q$ 在 Mellin 空间中与 $P^{(0)}_{\mathrm{NS}}$ 对易，于是 $P^{(1),\mathrm{DIS}}_{\mathrm{NS}}=P^{(1),\MSbar}_{\mathrm{NS}}-\frac{\beta_0}2C_2^q$。
（以上步骤均以符号计算核对。）
\end{proof}

\noindent
夸克分布的 DIS 方案可以追溯到\cite{AltarelliEllisMartinelli1979}；用动量求和规则固定胶子的约定参见\cite{DiemozEtAl1988}。其他过程（例如 Drell--Yan）的硬系数必须按同一个 $Z$ 变换，
在那里 DIS 方案的系数并不为零；方案的选择同时改变 PDF 与系数函数，可观测量只在所截断的阶以上改变。

% ---------------------------------------------------------------------
\section{重夸克质量与变味方案}\label{sec:disall-heavy}

粲与底夸克的质量 $m$ 介于 $\Lambda$ 与 $Q$ 之间：$Q\sim m$ 时重夸克只在末态成对产生，$Q\gg m$ 时 $\as\ln(Q^2/m^2)$ 变大，重夸克又应当作为部分子参与演化。
$\MSbar$ 与质量无关，重夸克不会自动退耦，退耦要通过 $n_f$ 与 $n_f-1$ 味有效理论的匹配来实现（\cref{rem:qcd-decoupling}）。下面 $Q$ 同时表示重夸克（如\cref{fig:disall-heavy}）与 $Q^2=-q^2$，从上下文可以区分。

\begin{theorem}[含重夸克质量的因子化（Collins 1998） \tagA]\label{thm:disall-heavy}
在硬部分中保留重夸克质量 $m_h$，并按 CWZ/ACOT 框架引入重夸克分布\cite{CollinsWilczekZee1978,AivazisEtAl1994}，DIS 结构函数仍取因子化形式
\begin{equation}\label{eq:disall-heavy}
\frac{F_a}{x}=\sum_i f_i\otimes\hat C_{a,i}\Big(\frac{m_h}Q,\frac Q\mu,\as(\mu)\Big)+R_a,
\end{equation}
其中 $i$ 包括重夸克，余项 $R_a$ 被 $\Lambda/Q$ 的某个幂次一致地压低，与 $m_h/Q$ 的大小无关：\cite{Collins1998b}的摘要写道
「“The power-suppressed corrections to the factorization formula are uniformly suppressed by a power of $\Lambda/Q$, independently of the size of heavy quark masses, $M$, relative to $Q$.”」
非微扰地位：\cref{asm:intro-bridge}（MA1、MA2）。
\uses{thm:disall-factorization,rem:qcd-decoupling,asm:intro-bridge}
% errtype: E-S
\end{theorem}

\begin{proofsketch}
论证与\cref{thm:disall-factorization}相同，只是重夸克质量保留在硬部分与相应的近似子中：$m_h\sim Q$ 时含重夸克的线是硬的；$m_h\ll Q$ 时重夸克也可以是活动部分子，其共线发散被 $m_h$ 截断为 $\ln(Q^2/m_h^2)$。
重整化采用 CWZ 的退耦减除方案\cite{CollinsWilczekZee1978}：$\mu<m_h$ 时理论等价于 $n_f-1$ 味的有效理论，$\mu>m_h$ 时用 $n_f$ 味的 $\MSbar$（\cref{rem:qcd-decoupling}）。
所得系数在 $Q\sim m_h$ 时趋于固定味方案（FFNS），在 $Q\gg m_h$ 时趋于零质量变味方案；一般质量变味方案（ACOT\cite{Aivazis:1993kh,AivazisEtAl1994}、S-ACOT\cite{Kramer:2000hn}、TR\cite{Thorne:1997ga}、
BMSN\cite{Buza:1996wv}、FONLL\cite{Forte:2010ta}）减去两者的重叠，在所计算的阶上是等价的处方（\cref{ex:disall-heavy-lo}）。完整论证见\cite{Collins1998b}。
\end{proofsketch}

\paragraph{地位与缺口}
本定理继承\cref{thm:disall-factorization}的全部瓶颈（\cref{asm:appN-collins,asm:appN-ward-product,asm:appN-cutsum}等）与具名缺口\cref{asm:appN-brst-nonlocal}；质量的保留不引入新的未证假设，余项的一致性同样是标度陈述（E-S）。
这一结论是对只有一个强子的 DIS 而言的：对强子—强子碰撞，\cite[第~14.6~节，第~569--570~页]{Collins2011}指出因子化只保证领头幂次上红外发散的相消，形如 $(m^2/Q^2)\times(1/\varepsilon)$ 的项可能残留，
在硬散射中保留入射重夸克的质量是“a bad idea”；把 ACOT 一类的方案移植到强子—强子碰撞时须记住这一区别。

\begin{example}[领头阶 $\gamma^*g\to Q\bar Q$]\label{ex:disall-heavy-lo}
考虑 $\gamma^*(q)+g(p)\to Q(p_1)+\bar Q(p_2)$，$p_1^2=p_2^2=m^2$（\cref{fig:disall-heavy}(a)）。记 $\epsilon_m\equiv m^2/Q^2$，$\hat s=Q^2(1-z)/z$，$\beta\equiv\sqrt{1-4m^2/\hat s}=\sqrt{1-4\epsilon_mz/(1-z)}$，阈值为 $z\le1/(1+4\epsilon_m)$。
与无质量情形相比，旋量求和改为 $\sla p_1+m$、$\sla p_2-m$，传播子分母为 $t_1\equiv(p-p_1)^2-m^2=-\frac{\hat s+Q^2}2(1-\beta\cos\theta)$ 与 $u_1\equiv(p-p_2)^2-m^2=-\frac{\hat s+Q^2}2(1+\beta\cos\theta)$（$\gamma^*g$ 质心系），
极化求和仍可取 $-g_{\rho\sigma}$（我们以物理极化求和核对过），四维两体相空间为 $\frac\beta{16\pi}\int_{-1}^1\dd\cos\theta$。用\cref{def:disall-breit}的投影（$D=4$）并对 $\cos\theta$ 积分，得到（本书的推导）
\begin{equation}\label{eq:disall-heavy-lo}
\begin{aligned}
\frac{\hat F^g_2}{ze_Q^2}&=\frac{\as}{2\pi}\,2T_F\Big\{\big[z^2+(1-z)^2+4\epsilon_mz(1-3z)-8\epsilon_m^2z^2\big]\ln\frac{1+\beta}{1-\beta}\\
&\hspace{8.2em}+\beta\big[8z(1-z)-1-4\epsilon_mz(1-z)\big]\Big\},\\
\frac{\hat F^g_L}{ze_Q^2}&=\frac{\as}{2\pi}\,2T_F\Big\{-8\epsilon_mz^2\ln\frac{1+\beta}{1-\beta}+4\beta z(1-z)\Big\}
\end{aligned}
\end{equation}
（阈值以下为零），与\cite[式~(49)--(53)]{Forte:2010ta}一致（那里 $a_s=\as/2\pi$，因子 2 同样计入 $Q$ 与 $\bar Q$）。对数 $\ln\frac{1+\beta}{1-\beta}$ 来自 $1/t_1$、$1/u_1$ 在 $\cos\theta\to\pm1$ 附近的增强：$m$ 正规化了重夸克与胶子共线的区域。
当 $m\to0$ 时 $\ln\frac{1+\beta}{1-\beta}=\ln\frac{Q^2}{m^2}+\ln\frac{1-z}z+\order{\epsilon_m}$，于是
$\hat F^g_2/(ze_Q^2)\to\frac{\as}{2\pi}2[P_{qg}\ln\frac{Q^2}{m^2}+C_2^g]$，$\hat F^g_L/(ze_Q^2)\to\frac{\as}{2\pi}2C_L^g$，修正为 $\order{\epsilon_m\ln\epsilon_m}$：
无质量计算中的极点 $-\frac2\varepsilon P_{qg}$ 被 $2P_{qg}\ln(Q^2/m^2)$ 代替，有限部分恰为 $\MSbar$ 的 $C_2^g$（\cref{prop:disall-nlo}）。另一方面，由\cref{def:pdf-quark}在光锥规范中直接计算单圈的重夸克—胶子分布
（$\bar Q$ 在壳、$Q$ 离壳的割图，$D$ 维横动量积分，$\MSbar$ 扣除紫外极点）得 $f_{Q/g}(z,\mu)=\frac{\as}{2\pi}P_{qg}(z)\ln\frac{\mu^2}{m^2}+\order{\as^2}$，没有附加常数。
于是\cref{fig:disall-heavy}(b) 的减除项（Born 取无质量形式，即 S-ACOT 的选择\cite{Kramer:2000hn}；$\mu<m$ 时按 CWZ 方案 $f_{Q/g}=0$，不作减除）为 $[f_{Q/g}+f_{\bar Q/g}]\otimes\hat F^{(0)}_2/z=e_Q^2\frac{\as}{2\pi}2P_{qg}\ln(\mu^2/m^2)$，减除后的胶子系数
$\hat C_{2,g}=\frac{\as}{2\pi}[2T_F\{\cdots\}-2P_{qg}\ln\frac{\mu^2}{m^2}]$（$\{\cdots\}$ 为\cref{eq:disall-heavy-lo}第一式的花括号）在 $\mu=Q$、$m\to0$ 时趋于 $\frac{\as}{2\pi}2C_2^g$；它与光滑函数的卷积对一切 $m/Q$ 一致有界（阈值以上未被抵消的 $\ln(Q^2/m^2)$ 只占宽度 $\order{\epsilon_m}$ 的区间）。
数值上，$z=0.1$ 时 $\{\cdots\}$ 与其无质量渐近式 $C_2^g/T_F+(P_{qg}/T_F)\ln(Q^2/m^2)$ 之差在 $\epsilon_m=0.1,\,0.01,\,0.001$ 时依次为 $0.074,\,0.014,\,0.002$。
在 $\order{\as}$，重夸克初态项 $e_Q^2(f_Q+f_{\bar Q})$ 加上减除后的胶子项，恰好回到固定味方案的结果：各方案在所计算的阶上等价，变味方案的好处是 $f_Q$ 的演化对 $\as\ln(\mu^2/m^2)$ 作了重求和。
\uses{thm:disall-heavy,def:disall-breit,prop:disall-nlo,def:pdf-quark}
\end{example}

%%% BEGIN fig:disall-heavy
\begin{figure}[htbp]
\centering
\begin{tikzpicture}[fd,
  fb lab/.style={font=\footnotesize,inner sep=1.5pt},
  fb heavy/.style={draw,solid,line width=1.5pt,postaction={decorate,solid,decoration={markings,
      mark=at position #1 with {\arrow[xshift=3.6pt,line width=0.8pt]{Stealth[length=7pt,width=6pt]}}}}},
  fb heavy/.default=0.5]
  % ======================= (a) LO photon--gluon fusion with massive quarks =======================
  \begin{scope}
    \coordinate (B) at (0,0.75); \coordinate (A) at (0,-0.75);
    \draw[photon] (-1.6,0.75) -- (B);
    \draw[gluon] (-1.6,-0.75) -- (A);
    \draw[fb heavy] (B) -- (1.55,0.75);
    \draw[fb heavy] (A) -- (B);
    \draw[fb heavy] (1.55,-0.75) -- (A);
    \node[vertex] at (A) {}; \node[vertex] at (B) {};
    \path[mom={$q$},mom sep=0.26cm] (-1.6,0.75) -- (B);
    \path[mom'={$k$},mom sep=0.3cm] (-1.6,-0.75) -- (A);
    \path[mom={$p_1$},mom sep=0.28cm] (B) -- (1.55,0.75);
    \path[mom'={$p_2$},mom sep=0.3cm] (A) -- (1.55,-0.75);
    \node[fb lab,anchor=east] at (-1.62,0.75) {$\gamma^*$};
    \node[fb lab,anchor=east] at (-1.62,-0.75) {$g$};
    \node[fb lab,anchor=west] at (1.57,0.75) {$Q$};
    \node[fb lab,anchor=west] at (1.57,-0.75) {$\bar Q$};
    \node[fb lab,anchor=west,align=left] at (2.3,0) {$+$\ 交叉图\\\phantom{$+$\ }（$p_1\leftrightarrow p_2$）};
    \node[font=\small,anchor=north] at (0.85,-1.45) {(a)\ 领头阶 $\gamma^*g\to Q\bar Q$（保留 $m$）};
  \end{scope}
  % ======================= (b) ACOT/CWZ subtraction term =======================
  \begin{scope}[xshift=8.15cm]
    \coordinate (B) at (0,0.75);
    \node[jet blob,ellipse,minimum width=1.55cm,minimum height=0.9cm,inner sep=0pt] (F) at (0,-0.75) {};
    \draw[gluon] (-1.75,-0.75) -- (F.west);
    \draw[photon] (-1.6,0.75) -- (B);
    \draw[fb heavy=0.55] (F.north) -- (B);
    \draw[fb heavy] (B) -- (1.55,0.75);
    \draw[fb heavy=0.6] (1.6,-1.25) -- (F.-25);
    \node[vertex] at (B) {};
    \node[font=\footnotesize] at (F) {$f_{Q/g}$};
    \path[mom={$q$},mom sep=0.26cm] (-1.6,0.75) -- (B);
    \path[mom'={$k$},mom sep=0.3cm] (-1.75,-0.75) -- (F.west);
    \node[fb lab,anchor=east] at (-1.62,0.75) {$\gamma^*$};
    \node[fb lab,anchor=east] at (-1.77,-0.75) {$g$};
    \node[fb lab,anchor=west] at (1.57,0.75) {$Q$};
    \node[fb lab,anchor=west] at (1.62,-1.25) {$\bar Q$};
    \node[fb lab,anchor=east] at (-0.06,0.02) {$Q$};
    \node[fb lab,anchor=west,text=gray!30!black] at (0.14,0.38) {$\hat\sigma(\gamma^*Q\to Q)$};
    \node[font=\small,anchor=north] at (0,-1.45) {(b)\ 减除项 $f_{Q/g}\otimes\hat\sigma(\gamma^*Q\to Q)$};
  \end{scope}
  \node[font=\Large] at (5.45,0) {$-$};
\end{tikzpicture}
\caption{含重夸克质量的 DIS（粗线为质量 $m$ 的重夸克 $Q$）。(a)~领头阶的光子—胶子融合，保留完整的质量依赖，两图之和的平方给出领头阶的胶子道系数；
(b)~变味（ACOT/CWZ）方案中的减除项 ${f_{Q/g}\otimes\hat\sigma(\gamma^*Q\to Q)}$：胶子先在重夸克分布 $f_{Q/g}$（蓝色，靶共线）中分裂为 $Q\bar Q$，其中的 $Q$ 再发生领头阶散射 $\gamma^*Q\to Q$。
二者之差在 $m\to0$ 时抵消 (a) 中的共线对数 $\ln(Q^2/m^2)P_{qg}$，给出对 $m/Q$ 一致有效的系数函数（\cref{thm:disall-heavy}，\cref{ex:disall-heavy-lo}）。}
\label{fig:disall-heavy}
\end{figure}
%%% END fig:disall-heavy

% ---------------------------------------------------------------------
\section{核靶}\label{sec:disall-nuclear}

原子核是色单态束缚态，DIS 的证明不用到靶的其他性质；核的特殊之处是线度随质量数 $A$ 增长，使某些幂次修正被增强，并在小 $x$ 处带来相干的多重散射。

\begin{proposition}[核靶上的 DIS \tagA]\label{prop:disall-nuclear}
设靶为质量数 $A$ 的原子核。
\begin{enumerate}[label=(\alph*),leftmargin=2em,itemsep=1pt,topsep=2pt]
\item 在领头幂次，核靶上的 DIS 与强子靶上的 DIS 具有相同的因子化（\cref{thm:disall-factorization}）：以核 PDF 代替强子 PDF，系数函数不变。
\item （核增强的扭度 4 修正，部分，\tagP）在高扭度的推广因子化中\cite{QiuSterman1991a}，核内的多重散射给出的扭度 4 修正被核的线度增强；
在 Luo--Qiu--Sterman 的处理中\cite{Luo:1993ui,Luo:1994np}，相应的扭度 4 核部分子关联函数按模型估计为 $\lambda^2A^{1/3}f_{i/A}$，$\lambda^2\sim0.05$--$0.1\,\mathrm{GeV}^2$，修正约为 $A^{1/3}\lambda^2/Q^2$。
\item 对强子—核碰撞（例如 pA 中的 Drell--Yan），初态相互作用可能使因子化形式失效，除非束流能量大于一个随靶长度增长的尺度（据\cite{Bodwin:1988fs}的摘要；较早的讨论见\cite{BodwinBrodskyLepage1981}）；
在只有一个强子、末态被求和的 DIS 中，这一初态机制不出现。
\item 小 $x$ 处的核遮蔽可以在领头扭度中，借助 Gribov--Glauber 理论与衍射 DIS 的因子化来描述\cite{Frankfurt:2011cs}。
\item （小 $x$ 的饱和区，猜想，\tagC）在足够小的 $x$ 与大核中，饱和效应可能使领头扭度的描述失效（\cref{ch:smallx}）。
\end{enumerate}
非微扰地位：\cref{asm:intro-bridge}（MA1、MA2）。
\uses{thm:disall-factorization,def:pdf-quark,def:pdf-gluon,asm:intro-bridge}
% errtype: E-S
\end{proposition}

\begin{proofsketch}
(a) 推导只用到靶是色单态的物理态；核 PDF 由\cref{def:pdf-quark,def:pdf-gluon}的同一组算符在核态中定义（通常按每个核子归一化），强子层次同样依赖于\cref{asm:intro-bridge}。
(b) 推广的高扭度因子化中\cite{QiuSterman1991a,QiuSterman1991b}，领头 $1/Q^2$ 修正是硬部分与扭度 2、扭度 4 矩阵元的卷积；$A^{1/3}$ 来自“两个场强位于同一核子之内”的模型假设（否则原则上可达 $A^{2/3}$）\cite{Guo:1997it}，
$\lambda^2$ 由光生双喷注的核依赖估计\cite{Luo:1993ui}。(c) DIS 中没有入射强子旁观者的初态相互作用，末态相互作用按\cref{thm:disall-factorization}的论证被求和。
(d) 遮蔽理论结合了 Gribov--Glauber 理论、QCD 因子化与 HERA 衍射 DIS 的分析\cite{Frankfurt:2011cs}（衍射 DIS 见\cref{ch:break}）。(e) 饱和区见\cref{ch:smallx}。
\end{proofsketch}

\paragraph{地位与缺口}
(a) 继承\cref{thm:disall-factorization}的全部瓶颈（\cref{asm:appN-collins,asm:appN-ward-product,asm:appN-cutsum}等），全局核 PDF 拟合都以领头扭度因子化为前提\cite{Eskola:2021nhw,Kovarik:2015cma,AbdulKhalek:2022fyi}。
真正未被控制的是幂次修正对核尺寸的依赖：(b) 只是模型估计（[P]），推广的高扭度因子化在核中没有全阶证明；(e) 是猜想（[C]），它与共线因子化的衔接没有被证明。

% ---------------------------------------------------------------------
\section{本章小结：对应表与常见误解}\label{sec:disall-summary}

本章的推理链是：非物理区的领头区域（\cref{prop:disall-two-routes}(a)）$\to$ 嵌套的减除项结构（\cref{prop:disall-factorized-subtraction}）$\to$ 与 $\MSbar$ 的对接（\cref{prop:disall-h5}；正规子路线\cref{thm:disall-renorm-compat}）
$\to$ 物理区：矩经由复圆盘上的 Cauchy 估计，固定 $x$ 经由割图分析（\cref{prop:disall-two-routes}(b)(c)）$\to$ 主定理\cref{thm:disall-factorization}（[A]），其强度恰好等于\cref{asm:appN-collins,asm:appN-ward-product,asm:appN-cutsum}等假设，
另加具名缺口\cref{asm:appN-brst-nonlocal}与\cref{asm:intro-bridge}。与附录~N 的对应见\cref{tab:disall-correspondence}。

% misconception box: M4 + critic O7; M25; M26
\begin{tcolorbox}[enhanced,breakable,title=本章涉及的常见误解,fonttitle=\bfseries,
  colback=gray!3!white,colframe=gray!55!black,boxrule=0.6pt,arc=2pt,
  left=5pt,right=5pt,top=3pt,bottom=3pt,before skip=6pt,after skip=6pt]
\begin{itemize}[leftmargin=1.2em,itemsep=1pt,topsep=1pt]
\item 「矩的因子化蕴含 $x$ 空间中一致的 $\order{\Lambda^2/Q^2}$ 余项」——矩只决定分布意义下的扭度 2 部分（\cref{cor:disall-consistency}）：对每个 $n$ 矩为 $\order{1/Q^2}$ 的 $\sin(Q^2x)$ 逐点为 $\order{1}$；
高扭度系数随 $n$ 增长，余项在 $x\to1$ 处不一致，共振区只能在平滑后比较。固定 $x$ 的逐点陈述是 [A] 定理\cref{thm:disall-factorization}，需要割图分析（\cref{prop:disall-two-routes}(c)）。
\item 「动量空间的 [A] 定理与 OPE 的 [R] 结果一致，所以逐点余项也是严格的」——一致的只是作为分布的扭度 2 部分（$R_{\mathrm{OPE}}$）；逐点余项只由全阶论证控制。
\item 「把 $\mu_R,\mu_F$ 变动 $2$ 倍就给出理论误差」——这是约定；定理只保证剩余的标度依赖高一阶（\cref{prop:disall-scale}）。
\item 「Callan--Gross 关系 $F_L=0$ 是 QCD 的定理」——它是部分子模型（领头阶）的结果；$F_L=\order{\as}$，$C_L^q=2C_Fz$，$C_L^g=4T_Fz(1-z)$（\cref{prop:disall-nlo}，Altarelli--Martinelli 关系）。
\end{itemize}
\end{tcolorbox}

\begin{table}[htbp]
\centering
\caption{本章的陈述与附录~N 中精确形式的对应。“误差类型”按\cref{def:intro-error-types}：E-S 为标度陈述，E-N 为积分范数界，E-N$^\varphi$ 为纤维化积分范数界。}
\label{tab:disall-correspondence}
\small
\begin{tabularx}{\textwidth}{@{}>{\raggedright\arraybackslash}p{0.2\textwidth}>{\raggedright\arraybackslash}p{0.26\textwidth}>{\raggedright\arraybackslash}X@{}}
\toprule
本章陈述 & 附录~N 中的精确形式 & 差别（所用假设；误差类型） \\
\midrule
\cref{prop:disall-two-routes}（$R_{\mathrm{cond}}$） & \cref{sec:appN-processes}；\cref{thm:appN-subtraction}、\cref{lem:appN-error-table} &
(a)：\cref{asm:appN-collins}、\cref{asm:appN-contour}、\cref{asm:appN-envelope}、\cref{asm:appN-separation}；(b) 另需复圆盘上的一致性；(c) 另需\cref{asm:appN-cutsum}。E-N。 \\
\cref{prop:disall-factorized-subtraction}（$R_{\mathrm{cond}}$） & \cref{sec:appN-processes}；\cref{prop:method-product}、\cref{lem:appN-algebra} &
\cref{asm:appN-ward-product}(WI1)(WI2)、\cref{asm:appN-uv}；被积函数层次的恒等式（不含估计）。 \\
\cref{thm:disall-renorm-compat}（$R_{\mathrm{cond}}$） & \cref{sec:appN-processes} &
\cref{asm:appN-multren}(Z1)(Z3)(i)、\cref{asm:appN-ir-regulator}(Z2)、\cref{asm:appN-contour}(b)；内容等级 [C]，不被任何 [A] 陈述使用。 \\
\cref{ex:disall-ladder} & \cref{prop:appN-ladder}（\cref{sec:appN-examples}） &
代数精确（$R_{\mathrm{PT}}$）；费米子两级梯形图的 (CS1)--(CS3) 未证明；标量模型单梯级梯形图的完整界见\cref{prop:appN-phi36-dis}（E-N$^\varphi$）。 \\
\cref{lem:disall-fsi}、\cref{prop:disall-regions}（[A]） & \cref{asm:appN-cutsum}；\cref{asm:appN-envelope}、\cref{asm:appN-separation} &
标度陈述（E-S）及其割图形式与积分形式。 \\
\cref{thm:disall-factorization}（[A]） & \cref{thm:appN-structure}（$R_{\mathrm{cond}}$） &
条件形式列出附录~N 的全部技术假设，包括具名缺口\cref{asm:appN-brst-nonlocal}，故其内容等级为 [C]；本章的 [A] 定理只点名该缺口。文献：E-S；条件形式：E-N。 \\
\bottomrule
\end{tabularx}
\end{table}

